% =============================================
% Объявление переменных (будут переопределены)
% =============================================
\newcommand{\funcCabinet}{}
\newcommand{\STO}{}
\newcommand{\units}{}
\newcommand{\functions}{}
\newcommand{\ManualUnitDefinite}{}

% =============================================
% Вспомогательные макросы для значений (для каждого РЭ)
% =============================================
\newcommand{\ManualOne}{%
  \renewcommand{\funcCabinet}{дифференциальной защиты для одной или двух секций сборных шин (не более 9 присоединений с жесткой фиксацией к шинам; контролируемых групп токов -- 11) на электрических станциях и подстанциях напряжением 110–220~кВ архитектуры I типа}%
  \renewcommand{\STO}{СТО 56947007-33.040.20.285-2019 ПАО «ФСК ЕЭС»}%
  \renewcommand{\units}{трёх интеллектуальных электронных устройств (ИЭУ) серии «ЮНИТ-М500» (для пофазного исполнения), имеющие код заказа ЮНИТ-М519-ДЗШ-Р102-К103-К103-К103-х-х-В101-В101-В101-В101-В101-М1С4.0С00-С01.03}%
  \renewcommand{\functions}{%
    \item дифференциальная защита шин (ДЗШ);
    \item контроль цепей тока (КЦТ);
    \item блокировка при неисправности цепей напряжения (БНН).}%
  \renewcommand{\ManualUnitDefinite}{ЮТКБ.656122.611-04.01 РЭ «Дифференциальная защита сборных шин для первичных схем до 12 элементов на напряжение от 110 до 220 кВ с функцией УРОВ. Устройство типа ЮНИТ-М500-ДЗШ»}%
}

\newcommand{\ManualTwo}{%
  \renewcommand{\funcCabinet}{дифференциальной защиты для двух секций сборных шин (не более 16 присоединений с жесткой фиксацией к шинам; контролируемых групп токов -- 19) на электрических станциях и подстанциях напряжением 110–220~кВ архитектуры I типа}%
  \renewcommand{\STO}{СТО 56947007-33.040.20.285-2019 ПАО «ФСК ЕЭС»}%
  \renewcommand{\units}{трёх интеллектуальных электронных устройств (ИЭУ) серии «ЮНИТ-М500» (для пофазного исполнения), имеющие код заказа ЮНИТ-М519-ДЗШ-Р102-К103-К103-К103-К103-К103-В121-В121-В101-М1С4.0С00-М1С4.0С00-С01.03}%
  \renewcommand{\functions}{%
    \item дифференциальная защита шин (ДЗШ);
    \item контроль цепей тока (КЦТ);
    \item блокировка при неисправности цепей напряжения (БНН).}%
  \renewcommand{\ManualUnitDefinite}{ЮТКБ.656122.611-04.02 РЭ «Дифференциальная защита сборных шин для первичных схем до 24 элементов на напряжение от 110 до 220 кВ с функцией УРОВ. Устройство типа ЮНИТ-М500-ДЗШ»}%
}

\newcommand{\ManualThree}{%
  \renewcommand{\funcCabinet}{дифференциальной защиты для двух секций сборных шин (не более 9 присоединений с управляемой фиксацией к шинам; контролируемых групп токов -- 11) на электрических станциях и подстанциях напряжением 110–220~кВ архитектуры I типа}%
  \renewcommand{\STO}{СТО 56947007-33.040.20.285-2019 ПАО «ФСК ЕЭС»}%
  \renewcommand{\units}{трёх интеллектуальных электронных устройств (ИЭУ) серии «ЮНИТ-М500» (для пофазного исполнения), имеющие код заказа ЮНИТ-М519-ДЗШ-Р102-К103-К103-К103-х-х-В101-В101-В101-В101-В101-М1С4.0С00-С01.03}%
  \renewcommand{\functions}{%
    \item дифференциальная защита шин (ДЗШ);
    \item контроль цепей тока (КЦТ);
    \item блокировка при неисправности цепей напряжения (БНН).}%
  \renewcommand{\ManualUnitDefinite}{ЮТКБ.656122.611-04.01 РЭ «Дифференциальная защита сборных шин для первичных схем до 12 элементов на напряжение от 110 до 220 кВ с функцией УРОВ. Устройство типа ЮНИТ-М500-ДЗШ»}%
}

\newcommand{\ManualFour}{%
  \renewcommand{\funcCabinet}{дифференциальной защиты для двух секций сборных шин (не более 15 присоединений с управляемой фиксацией к шинам; контролируемых групп токов -- 20) на электрических станциях и подстанциях напряжением 110–220~кВ архитектуры I типа}%
  \renewcommand{\STO}{СТО 56947007-33.040.20.285-2019 ПАО «ФСК ЕЭС»}%
  \renewcommand{\units}{трёх интеллектуальных электронных устройств (ИЭУ) серии «ЮНИТ-М500» (для пофазного исполнения), имеющие код заказа ЮНИТ-М519-ДЗШ-Р102-К103-К103-К103-В121-В121-В121-В121-В101-М1С4.0С00-М1С4.0С00-С01.03}%
  \renewcommand{\functions}{%
    \item дифференциальная защита шин (ДЗШ);
    \item контроль цепей тока (КЦТ);
    \item блокировка при неисправности цепей напряжения (БНН).}%
  \renewcommand{\ManualUnitDefinite}{ЮТКБ.656122.611-04.02 РЭ «Дифференциальная защита сборных шин для первичных схем до 24 элементов на напряжение от 110 до 220 кВ с функцией УРОВ. Устройство типа ЮНИТ-М500-ДЗШ»}%
}

\newcommand{\ManualFive}{%
  \renewcommand{\funcCabinet}{защиты ошиновок ВН, СН, НН (авто)трансформаторов и ошиновок оборудования СКРМ на электрических станциях и подстанциях напряжением 6–750~кВ архитектуры I типа}%
  \renewcommand{\STO}{СТО 56947007-33.040.20.285-2019 ПАО «ФСК ЕЭС»}%
  \renewcommand{\units}{двух интеллектуальных электронных устройств (ИЭУ) серии «ЮНИТ-М500» (два комплекта), имеющие код заказа ЮНИТ-М514-ДЗО-Р102-К103-К103-К103-х-В101-В101-М1С4.0С00-С01.03}%
  \renewcommand{\functions}{%
    \item дифференциальная защита ошиновки (ДЗО);
    \item контроль цепей тока (КЦТ);
    \item блокировка при неисправности цепей напряжения (БНН).}%
  \renewcommand{\ManualUnitDefinite}{ЮТКБ.656122.611-05.01 РЭ «Дифференциальная защита ошиновки от 6 до 750~кВ. Устройство типа ЮНИТ-М500-ДЗО»}%
}

\newcommand{\ManualSix}{%
  \renewcommand{\funcCabinet}{защиты ошиновок ВН, СН, НН (авто)трансформаторов и ошиновок оборудования СКРМ на электрических станциях и подстанциях напряжением 6–750~кВ архитектуры I типа}%
  \renewcommand{\STO}{СТО 56947007-33.040.20.285-2019 ПАО «ФСК ЕЭС»}%
  \renewcommand{\units}{одного интеллектуального электронного устройств (ИЭУ) серии «ЮНИТ-М500», имеющего код заказа ЮНИТ-М514-ДЗО-Р102-К103-К103-К103-х-В101-В101-М1С4.0С00-С01.03}%
  \renewcommand{\functions}{%
    \item дифференциальная защита ошиновки (ДЗО);
    \item контроль цепей тока (КЦТ);
    \item блокировка при неисправности цепей напряжения (БНН).}%
  \renewcommand{\ManualUnitDefinite}{ЮТКБ.656122.611-05.01 РЭ «Дифференциальная защита ошиновки от 6 до 750~кВ. Устройство типа ЮНИТ-М500-ДЗО»}%
}

\newcommand{\ManualSeven}{%
  \renewcommand{\funcCabinet}{ступенчатых защит, УРОВ и АУВ обходного выключателя (ОВ) на электрических станциях и подстанциях напряжением 110–220~кВ архитектуры I типа}%
  \renewcommand{\STO}{СТО 56947007-33.040.20.285-2019 ПАО «ФСК ЕЭС»}%
  \renewcommand{\units}{одного интеллектуального электронного устройств (ИЭУ) серии «ЮНИТ-М500», имеющего код заказа ЮНИТ-М514-ЛВ-Р102-К103-К103-К103-В101-В101-В101-М179.0700-С01.03}%
  \renewcommand{\functions}{%
    \item дистанционная защита (ДЗ);
	\item токовая защита нулевой последовательности (ТЗНП);
	\item максимальная токовая защита (МТЗ);
	\item междуфазная токовая отсечка (МФТО);
	\item оперативное ускорение (ОУ);
	\item автоматическое ускорение (АУ);
	\item блокировка при неисправности цепей напряжения (БНН);
	\item блокировка при качаниях (БК);
	\item защита от неполнофазного режима (ЗНР);
	\item определение места повреждения (ОМП);
	\item устройство резервирования отказа выключателя (УРОВ);
	\item автоматическое повторное включение (АПВ);
	\item автоматика управления выключателем (АУВ).}%
  \renewcommand{\ManualUnitDefinite}{ЮТКБ.656122.611-03.01 РЭ «Комплект ступенчатых защит для линии напряжением от 110 до 220~кВ с функциями УРОВ и АУВ. Устройство типа ЮНИТ-М500-ЛВ»}%
}

\newcommand{\ManualEight}{%
  \renewcommand{\funcCabinet}{ступенчатых защит, УРОВ и АУВ шиносоединительного (секционного) выключателя на электрических станциях и подстанциях напряжением 110–750~кВ архитектуры I типа}%
  \renewcommand{\STO}{СТО 56947007-33.040.20.285-2019 ПАО «ФСК ЕЭС»}%
  \renewcommand{\units}{одного интеллектуального электронного устройств (ИЭУ) серии «ЮНИТ-М500», имеющего код заказа ЮНИТ-М514-СВ-Р102-К103-К103-х-В101-В101-В101-М139.0300-С01.03}%
  \renewcommand{\functions}{%
    \item дистанционная защита (ДЗ);
	\item токовая защита нулевой последовательности (ТЗНП);
	\item максимальная токовая защита (МТЗ);
	\item междуфазная токовая отсечка (МФТО);
	\item оперативное ускорение (ОУ);
	\item автоматическое ускорение (АУ);
	\item блокировка при неисправности цепей напряжения (БНН);
	\item блокировка при качаниях (БК);
	\item защита от неполнофазного режима (ЗНР);
	\item автоматическое опережающее деление сети (АОДС);
	\item устройство резервирования отказа выключателя (УРОВ);
	\item автоматическое повторное включение (АПВ);
	\item контроль синхронизма и напряжения (КСН);
	\item защита от непереключения фаз (ЗНФ);
	\item автоматика управления выключателем (АУВ).}%
  \renewcommand{\ManualUnitDefinite}{ЮТКБ.656122.611-08.01 РЭ «Комплект защит шиносоединительного/секционного выключателя с функциями УРОВ и АУВ. Устройство типа ЮНИТ-М500-СВ»}%
}

\newcommand{\ManualNine}{%
  \renewcommand{\funcCabinet}{автоматики управления выключателем (АУВ) и устройства резервирования отказа выключателя (УРОВ) на электрических станциях и подстанциях напряжением 110–220~кВ архитектуры I типа}%
  \renewcommand{\STO}{СТО 56947007-33.040.20.285-2019 ПАО «ФСК ЕЭС»}%
  \renewcommand{\units}{одного интеллектуального электронного устройства (ИЭУ) серии «ЮНИТ-М500», имеющего код заказа ЮНИТ-М514-АУВ-Р102-К103-К103-К103-В101-В101-В101-М179.0700-С01.03}%
  \renewcommand{\functions}{%
    \item контроль синхронизма и напряжений (КСН);
    \item автоматическое повторное включение (АПВ);
    \item полуавтоматическое включение (ПАВ);
    \item контроль силового выключателя (КСВ);
    \item управление выключателем (УВ);
    \item фиксация положения выключателя (ФПВ);
    \item фиксация состояния ЛЭП (ФСЛ);
    \item устройство резервирования отказа выключателя (УРОВ).}%
  \renewcommand{\ManualUnitDefinite}{ЮТКБ.656122.611-10.01 РЭ «Автоматика управления и УРОВ выключателя напряжением от 110 до 220 кВ. Устройство типа ЮНИТ-М500-АУВ»}%
}

\newcommand{\ManualTen}{%
  \renewcommand{\funcCabinet}{автоматики управления выключателем (АУВ) и устройства резервирования отказа выключателя (УРОВ) СКРМ на электрических станциях и подстанциях напряжением 110–220~кВ архитектуры I типа}%
  \renewcommand{\STO}{СТО 56947007-33.040.20.285-2019 ПАО «ФСК ЕЭС»}%
  \renewcommand{\units}{одного интеллектуального электронного устройства (ИЭУ) серии «ЮНИТ-М500», имеющего код заказа ЮНИТ-М519-АУВ-Р102-К103-К103-К103-х-х-В101-В101-В101-В101-В101-М141.0400-С01.03}%
  \renewcommand{\functions}{%
    \item устройство резервирования отказа выключателя (УРОВ);
	\item защита от непереключения фаз (ЗНФ);
    \item автоматика управления выключателем (АУВ).}%
  \renewcommand{\ManualUnitDefinite}{ЮТКБ.656122.611-10.02 РЭ «Автоматика управления и УРОВ выключателя СКРМ напряжением от 110 до 220 кВ. Устройство типа ЮНИТ-М500-АУВ»}%
}

\newcommand{\ManualEleven}{%
  \renewcommand{\funcCabinet}{дифференциально-фазной защиты ЛЭП с комплектом ступенчатых защит и обменом разрешающими сигналами (для схем РУ с одним выключателем или с ОСШ) на электрических станциях и подстанциях напряжением 110–220~кВ архитектуры I типа}%
  \renewcommand{\STO}{СТО 56947007-33.040.20.282-2019 ПАО «ФСК ЕЭС»}%
  \renewcommand{\units}{одного интеллектуального электронного устройства (ИЭУ) серии «ЮНИТ-М500», имеющего код заказа ЮНИТ-М514-ВЧ-Р102-К103-К103-В101-В101-В120-М179.0700-В101-С01.03}%
  \renewcommand{\functions}{%
    \item дифференциально-фазная защита (ДФЗ);
	\item дистанционная защита (ДЗ);
	\item токовая защита нулевой последовательности (ТЗНП);
	\item максимальная токовая защита (МТЗ);
	\item междуфазная токовая отсечка (МФТО);
	\item оперативное ускорение (ОУ);
	\item автоматическое ускорение (АУ);
	\item блокировка при неисправности цепей напряжения (БНН);
	\item блокировка при качаниях (БК);
	\item ЛС (осуществляет передачу дискретных сигналов ТУ/ТО по ВЧ каналу с использованием возможностей ВЧ приемопередатчика);
	\item защита от неполнофазного режима (ЗНР);
    \item определение места повреждения (ОМП).}%
  \renewcommand{\ManualUnitDefinite}{ЮТКБ.656122.611-01.01 РЭ «Высокочастотные защиты с комплектом ступенчатых защит для линии напряжением от 110 до 220~кВ с функциями УРОВ и АУВ. Устройство типа ЮНИТ-М500-ВЧ»}%
}

\newcommand{\ManualTwelve}{%
  \renewcommand{\funcCabinet}{дифференциально-фазной защиты ЛЭП с комплектом ступенчатых защит и обменом разрешающими сигналами (для ЛЭП, подключенных через два выключателя или через один выключатель с возможностью перевода на ОСШ) на электрических станциях и подстанциях напряжением 110–220~кВ архитектуры I типа}%
  \renewcommand{\STO}{СТО 56947007-33.040.20.282-2019 ПАО «ФСК ЕЭС»}%
  \renewcommand{\units}{одного интеллектуального электронного устройства (ИЭУ) серии «ЮНИТ-М500», имеющего код заказа ЮНИТ-М514-ВЧ-Р102-К103-К103-В101-В101-В120-М179.0700-В101-С01.03}%
  \renewcommand{\functions}{%
    \item дифференциально-фазная защита (ДФЗ);
	\item дистанционная защита (ДЗ);
	\item токовая защита нулевой последовательности (ТЗНП);
	\item максимальная токовая защита (МТЗ);
	\item междуфазная токовая отсечка (МФТО);
	\item оперативное ускорение (ОУ);
	\item автоматическое ускорение (АУ);
	\item блокировка при неисправности цепей напряжения (БНН);
	\item блокировка при качаниях (БК);
	\item ЛС (осуществляет передачу дискретных сигналов ТУ/ТО по ВЧ каналу с использованием возможностей ВЧ приемопередатчика);
	\item защита от неполнофазного режима (ЗНР);
    \item определение места повреждения (ОМП).}%
  \renewcommand{\ManualUnitDefinite}{ЮТКБ.656122.611-01.01 РЭ «Высокочастотные защиты с комплектом ступенчатых защит для линии напряжением от 110 до 220~кВ с функциями УРОВ и АУВ. Устройство типа ЮНИТ-М500-ВЧ»}%
}

\newcommand{\ManualThirteen}{%
  \renewcommand{\funcCabinet}{направленной высокочастотной защиты ЛЭП с комплектом ступенчатых защит и обменом разрешающими сигналами (для ЛЭП, подключенных через один выключатель без возможности перевода на ОСШ) на электрических станциях и подстанциях напряжением 110–220~кВ архитектуры I типа}%
  \renewcommand{\STO}{СТО 56947007-33.040.20.282-2019 ПАО «ФСК ЕЭС»}%
  \renewcommand{\units}{одного интеллектуального электронного устройства (ИЭУ) серии «ЮНИТ-М500», имеющего код заказа ЮНИТ-М514-ВЧ-Р102-К103-К103-В101-В101-В120-М179.0700-В101-С01.03}%
  \renewcommand{\functions}{%
    \item направленная высокочастотная защита (НВЧЗ);
	\item дистанционная защита (ДЗ);
	\item токовая защита нулевой последовательности (ТЗНП);
	\item максимальная токовая защита (МТЗ);
	\item междуфазная токовая отсечка (МФТО);
	\item оперативное ускорение (ОУ);
	\item автоматическое ускорение (АУ);
	\item блокировка при неисправности цепей напряжения (БНН);
	\item блокировка при качаниях (БК);
	\item ЛС (осуществляет передачу дискретных сигналов ТУ/ТО по ВЧ каналу с использованием возможностей ВЧ приемопередатчика);
	\item защита от неполнофазного режима (ЗНР);
    \item определение места повреждения (ОМП).}%
  \renewcommand{\ManualUnitDefinite}{ЮТКБ.656122.611-01.01 РЭ «Высокочастотные защиты с комплектом ступенчатых защит для линии напряжением от 110 до 220~кВ с функциями УРОВ и АУВ. Устройство типа ЮНИТ-М500-ВЧ»}%
}

\newcommand{\ManualFourteen}{%
  \renewcommand{\funcCabinet}{направленной высокочастотной защиты ЛЭП с комплектом ступенчатых защит и обменом разрешающими сигналами (для ЛЭП, подключенных через два выключателя или через один выключатель с возможностью перевода на ОСШ) на электрических станциях и подстанциях напряжением 110–220~кВ архитектуры I типа}%
  \renewcommand{\STO}{СТО 56947007-33.040.20.282-2019 ПАО «ФСК ЕЭС»}%
  \renewcommand{\units}{одного интеллектуального электронного устройства (ИЭУ) серии «ЮНИТ-М500», имеющего код заказа ЮНИТ-М514-ВЧ-Р102-К103-К103-В101-В101-В120-М179.0700-В101-С01.03}%
  \renewcommand{\functions}{%
    \item направленная высокочастотная защита (НВЧЗ);
	\item дистанционная защита (ДЗ);
	\item токовая защита нулевой последовательности (ТЗНП);
	\item максимальная токовая защита (МТЗ);
	\item междуфазная токовая отсечка (МФТО);
	\item оперативное ускорение (ОУ);
	\item автоматическое ускорение (АУ);
	\item блокировка при неисправности цепей напряжения (БНН);
	\item блокировка при качаниях (БК);
	\item ЛС (осуществляет передачу дискретных сигналов ТУ/ТО по ВЧ каналу с использованием возможностей ВЧ приемопередатчика);
	\item защита от неполнофазного режима (ЗНР);
    \item определение места повреждения (ОМП).}%
  \renewcommand{\ManualUnitDefinite}{ЮТКБ.656122.611-01.01 РЭ «Высокочастотные защиты с комплектом ступенчатых защит для линии напряжением от 110 до 220~кВ с функциями УРОВ и АУВ. Устройство типа ЮНИТ-М500-ВЧ»}%
}

\newcommand{\ManualFiveteen}{%
  \renewcommand{\funcCabinet}{дифференциальной защиты ЛЭП с комплектом ступенчатых защит и обменом разрешающими сигналами (для ЛЭП, подключенных через один выключатель без возможности перевода на ОСШ) на электрических станциях и подстанциях напряжением 110–220~кВ архитектуры I типа}%
  \renewcommand{\STO}{СТО 56947007-33.040.20.282-2019 ПАО «ФСК ЕЭС»}%
  \renewcommand{\units}{одного интеллектуального электронного устройства (ИЭУ) серии «ЮНИТ-М500», имеющего код заказа ЮНИТ-М510-ДЗЛ-Р102-К103-В101-В101-М179.0700-С04.03}%
  \renewcommand{\functions}{%
    \item дифференциальная защита линии (ДЗЛ);
	\item дистанционная защита (ДЗ);
	\item токовая защита нулевой последовательности (ТЗНП);
	\item максимальная токовая защита (МТЗ);
	\item междуфазная токовая отсечка (МФТО);
	\item оперативное ускорение (ОУ);
	\item автоматическое ускорение (АУ);
	\item блокировка при неисправности цепей напряжения (БНН);
	\item блокировка при качаниях (БК);
	\item ЛС (осуществляет передачу дискретных сигналов ТУ/ТО по ВЧ каналу с использованием возможностей ВЧ приемопередатчика);
	\item защита от неполнофазного режима (ЗНР);
    \item определение места повреждения (ОМП).}%
  \renewcommand{\ManualUnitDefinite}{ЮТКБ.656122.611-02.01 РЭ «Дифференциальная защита линии с комплектом ступенчатых защит для линии напряжением от 110 до 220~кВ с функциями УРОВ и АУВ. Устройство типа ЮНИТ-М500-ДЗЛ»}%
}

\newcommand{\ManualSixteen}{%
  \renewcommand{\funcCabinet}{дифференциальной защиты ЛЭП с комплектом ступенчатых защит и обменом разрешающими сигналами (для ЛЭП, подключенных через два выключателя или через один выключатель с возможностью перевода на ОСШ) на электрических станциях и подстанциях напряжением 110–220~кВ архитектуры I типа}%
  \renewcommand{\STO}{СТО 56947007-33.040.20.282-2019 ПАО «ФСК ЕЭС»}%
  \renewcommand{\units}{одного интеллектуального электронного устройства (ИЭУ) серии «ЮНИТ-М500», имеющего код заказа ЮНИТ-М510-ДЗЛ-Р102-К103-В101-В101-М179.0700-С04.03}%
  \renewcommand{\functions}{%
    \item дифференциальная защита линии (ДЗЛ);
	\item дистанционная защита (ДЗ);
	\item токовая защита нулевой последовательности (ТЗНП);
	\item максимальная токовая защита (МТЗ);
	\item междуфазная токовая отсечка (МФТО);
	\item оперативное ускорение (ОУ);
	\item автоматическое ускорение (АУ);
	\item блокировка при неисправности цепей напряжения (БНН);
	\item блокировка при качаниях (БК);
	\item ЛС (осуществляет передачу дискретных сигналов ТУ/ТО по ВЧ каналу с использованием возможностей ВЧ приемопередатчика);
	\item защита от неполнофазного режима (ЗНР);
    \item определение места повреждения (ОМП).}%
  \renewcommand{\ManualUnitDefinite}{ЮТКБ.656122.611-02.01 РЭ «Дифференциальная защита линии с комплектом ступенчатых защит для линии напряжением от 110 до 220~кВ с функциями УРОВ и АУВ. Устройство типа ЮНИТ-М500-ДЗЛ»}%
}

\newcommand{\ManualSeventeen}{%
  \renewcommand{\funcCabinet}{ступенчатых защит с обменом разрешающими сигналами (для ЛЭП, подключенных через один выключатель без возможности перевода на ОСШ) на электрических станциях и подстанциях напряжением 110–220~кВ архитектуры I типа}%
  \renewcommand{\STO}{СТО 56947007-33.040.20.282-2019 ПАО «ФСК ЕЭС»}%
  \renewcommand{\units}{одного интеллектуального электронного устройства (ИЭУ) серии «ЮНИТ-М500», имеющего код заказа ЮНИТ-М514-ЛВ-Р102-К103-К103-К103-В101-В101-В101-М179.0700-С01.03}%
  \renewcommand{\functions}{%
	\item дистанционная защита (ДЗ);
	\item токовая защита нулевой последовательности (ТЗНП);
	\item максимальная токовая защита (МТЗ);
	\item междуфазная токовая отсечка (МФТО);
	\item оперативное ускорение (ОУ);
	\item автоматическое ускорение (АУ);
	\item блокировка при неисправности цепей напряжения (БНН);
	\item блокировка при качаниях (БК);
	\item логика связи (ЛС);
	\item защита от неполнофазного режима (ЗНР);
    \item определение места повреждения (ОМП).}%
  \renewcommand{\ManualUnitDefinite}{ЮТКБ.656122.611-03.01 РЭ «Комплект ступенчатых защит для линии напряжением от 110 до 220~кВ с функциями УРОВ и АУВ. Устройство типа ЮНИТ-М500-ЛВ»}%
}

\newcommand{\ManualEighteen}{%
  \renewcommand{\funcCabinet}{ступенчатых защит с обменом разрешающими сигналами (для ЛЭП, подключенных через два выключателя или через один выключатель с возможностью перевода на ОСШ) на электрических станциях и подстанциях напряжением 110–220~кВ архитектуры I типа}%
  \renewcommand{\STO}{СТО 56947007-33.040.20.282-2019 ПАО «ФСК ЕЭС»}%
  \renewcommand{\units}{одного интеллектуального электронного устройства (ИЭУ) серии «ЮНИТ-М500», имеющего код заказа ЮНИТ-М514-ЛВ-Р102-К103-К103-К103-В101-В101-В101-М179.0700-С01.03}%
  \renewcommand{\functions}{%
	\item дистанционная защита (ДЗ);
	\item токовая защита нулевой последовательности (ТЗНП);
	\item максимальная токовая защита (МТЗ);
	\item междуфазная токовая отсечка (МФТО);
	\item оперативное ускорение (ОУ);
	\item автоматическое ускорение (АУ);
	\item блокировка при неисправности цепей напряжения (БНН);
	\item блокировка при качаниях (БК);
	\item логика связи (ЛС);
	\item защита от неполнофазного режима (ЗНР);
    \item определение места повреждения (ОМП).}%
  \renewcommand{\ManualUnitDefinite}{ЮТКБ.656122.611-03.01 РЭ «Комплект ступенчатых защит для линии напряжением от 110 до 220~кВ с функциями УРОВ и АУВ. Устройство типа ЮНИТ-М500-ЛВ»}%
}

\newcommand{\ManualNineteen}{%
  \renewcommand{\funcCabinet}{ступенчатых защит с обменом разрешающими сигналами, УРОВ, АУВ (для ЛЭП, подключенных через один выключатель без возможности перевода на ОСШ) на электрических станциях и подстанциях напряжением 110–220~кВ архитектуры I типа}%
  \renewcommand{\STO}{СТО 56947007-33.040.20.282-2019 ПАО «ФСК ЕЭС»}%
  \renewcommand{\units}{одного интеллектуального электронного устройства (ИЭУ) серии «ЮНИТ-М500», имеющего код заказа ЮНИТ-М514-ЛВ-Р102-К103-К103-К103-В101-В101-В101-М179.0700-С01.03}%
  \renewcommand{\functions}{%
	\item дистанционная защита (ДЗ);
	\item токовая защита нулевой последовательности (ТЗНП);
	\item максимальная токовая защита (МТЗ);
	\item междуфазная токовая отсечка (МФТО);
	\item оперативное ускорение (ОУ);
	\item автоматическое ускорение (АУ);
	\item блокировка при неисправности цепей напряжения (БНН);
	\item блокировка при качаниях (БК);
	\item логика связи (ЛС);
	\item защита от неполнофазного режима (ЗНР);
    \item определение места повреждения (ОМП).}%
  \renewcommand{\ManualUnitDefinite}{ЮТКБ.656122.611-03.01 РЭ «Комплект ступенчатых защит для линии напряжением от 110 до 220~кВ с функциями УРОВ и АУВ. Устройство типа ЮНИТ-М500-ЛВ»}%
}

\newcommand{\ManualTwenty}{%
  \renewcommand{\funcCabinet}{защит трансформатора 110–220~кВ и ошиновки 6–35~кВ на электрических станциях и подстанциях архитектуры I типа}%
  \renewcommand{\STO}{СТО 56947007-33.040.20.276-2019 ПАО «ФСК ЕЭС»}%
  \renewcommand{\units}{двух интеллектуальных электронных устройств (ИЭУ) серии «ЮНИТ-М500», имеющие коды заказа ЮНИТ-М519-ДЗТ-Р102-К103-К103-К103-В101-В101-В101-В101-В121-М186.0800-М186.0800-С01.03 и ЮНИТ-М514-ДЗО-Р102-К103-К103-К103-х-В101-В101-М1С4.0С00-С01.03}%
  \renewcommand{\functions}{%
    \item дифференциальная защита трансформатора (ДЗТ);
    \item газовая защита (ГЗ);
    \item технологические защиты (ТЗ);
    \item устройство резервирования отказа выключателя НН (ПО УРОВ НН);
    \item максимальная токовая защита ВН (МТЗ ВН);
    \item максимальная токовая защита СН(НН1) (МТЗ СН(НН1));
    \item максимальная токовая защита СН(НН2) (МТЗ СН(НН2));
    \item комбинированный пуск по напряжению (НН, СН)1--(НН, СН)4 (КПН (НН, СН)1--(НН, СН)4);
    \item токовая защита нулевой последовательности ВН (ТЗНП ВН);
    \item токовая защита нулевой последовательности РЗН1, РЗН2 (ТЗНП РЗН1, РЗН2);
    \item защита от неполнофазного режима (ЗНР);
    \item защита от перегрузки (ЗП);
    \item контроль изоляции стороны НН(СН) (КИ НН(СН));
    \item контроль отсутствия напряжения трансформатора (КОН Т);
    \item токовый орган блокировки РПН (ТО блок. РПН);
    \item реле тока пуска охлаждения (РТ ПО);
    \item токовый контроль для защиты от дуговых замыканий (ТК ЗДЗ);
    \item блокировка при неисправности цепей напряжения сторон НН(СН);
    \item дифференциальная защита ошиновки НН(СН) (ДЗО НН(СН)).}%
  \renewcommand{\ManualUnitDefinite}{ЮТКБ.656122.611-06.01 РЭ «Основная защита трансформатора напряжением от 110 до 220~кВ с функциями УРОВ и АУВ. Устройство типа ЮНИТ-М500-ДЗТ» и ЮТКБ.656122.611-05.01 РЭ «Дифференциальная защита ошиновки от 6 до 750~кВ. Устройство типа ЮНИТ-М500-ДЗО»}%
}

\newcommand{\ManualTwentyOne}{%
  \renewcommand{\funcCabinet}{защит трансформатора 110–220~кВ на электрических станциях и подстанциях архитектуры I типа}%
  \renewcommand{\STO}{СТО 56947007-33.040.20.276-2019 ПАО «ФСК ЕЭС»}%
  \renewcommand{\units}{одного интеллектуального электронного устройства (ИЭУ) серии «ЮНИТ-М500», имеющего код заказа ЮНИТ-М519-ДЗТ-Р102-К103-К103-К103-В101-В101-В101-В101-В121-М186.0800-М186.0800-С01.03}%
  \renewcommand{\functions}{%
    \item дифференциальная защита трансформатора (ДЗТ);
    \item газовая защита трансформатора (ГЗ Т);
    \item технологические защиты трансформатора (ТЗ Т);
	\item газовая защита устройство регулирования напряжения под нагрузкой (ГЗ РПН);
    \item контроль изоляции газовых защит (КИ ГЗ);
	\item контроль изоляции технологических защит (КИ ТЗ);
	\item максимальная токовая защита с компбинированным пуском по напряжению стороны ВН (МТЗ/U ВН);
    \item токовая защита нулевой последовательности ВН (ТЗНП ВН);
	\item максимальная токовая защита с компбинированным пуском по напряжению стороны НН(СН) (МТЗ/U (НН, СН));
	\item дифференциальная защита ошиновки НН(СН) (ДЗО НН(СН));
	\item защита от неполнофазного режима (ЗНР);
    \item токовая защита нулевой последовательности РЗН1, РЗН2 (ТЗНП РЗН1, РЗН2);	
    \item контроль изоляции стороны НН(СН) (КИ НН(СН));
	\item блокировка при неисправности цепей напряжения сторон НН(СН) (БНН НН(СН));
	\item защита от перегрузки стороны ВН (ЗП ВН);
	\item защита от перегрузки стороны НН(СН) (ЗП НН(СН));
	\item контроль отсутствия напряжения трансформатора (КОН Т);
	\item токовый контроль для защиты от дуговых замыканий (ТК ЗДЗ);
	\item устройство резервирования отказа выключателя НН (ПО УРОВ НН);
	\item реле тока пуска охлаждения (РТ ПО).}%
  \renewcommand{\ManualUnitDefinite}{ЮТКБ.656122.611-06.01 РЭ «Основная защита трансформатора напряжением от 110 до 220~кВ с функциями УРОВ и АУВ. Устройство типа ЮНИТ-М500-ДЗТ»}%
}

\newcommand{\ManualTwentyTwo}{%
  \renewcommand{\funcCabinet}{резервных защит трансформатора 110–220~кВ на электрических станциях и подстанциях архитектуры I типа}%
  \renewcommand{\STO}{СТО 56947007-33.040.20.276-2019 ПАО «ФСК ЕЭС»}%
  \renewcommand{\units}{одного интеллектуального электронного устройства (ИЭУ) серии «ЮНИТ-М500», имеющего код заказа ЮНИТ-М514-РЗТ-Р102-К103-К103-х-В101-В101-В101-М168.0600-С01.03}%
  \renewcommand{\functions}{%
    \item газовая защита трансформатора (ГЗ Т);
	\item газовая защита устройство регулирования напряжения под нагрузкой (ГЗ РПН);
    \item контроль изоляции газовых защит (КИ ГЗ);
	\item максимальная токовая защита с компбинированным пуском по напряжению стороны ВН (МТЗ/U ВН);
    \item токовая защита нулевой последовательности ВН (ТЗНП ВН);
	\item максимальная токовая защита с компбинированным пуском по напряжению стороны СН (МТЗ/U СН);
	\item защита от неполнофазного режима (ЗНР);
	\item блокировка при неисправности цепей напряжения стороны НН(СН) (БНН НН(СН)).}%
  \renewcommand{\ManualUnitDefinite}{ЮТКБ.656122.611-07.01 РЭ «Резервные защиты трансформатора напряжением от 110 до 220 кВ. Устройство типа ЮНИТ-М500-РЗТ»}%
}

\newcommand{\ManualTwentyThree}{%
  \renewcommand{\funcCabinet}{основных защит автотрансформатора 220~кВ и ошиновки 6-35~кВ на электрических станциях и подстанциях архитектуры I типа}%
  \renewcommand{\STO}{СТО 56947007-33.040.20.276-2019 ПАО «ФСК ЕЭС»}%
  \renewcommand{\units}{двух интеллектуальных электронных устройств (ИЭУ) серии «ЮНИТ-М500», имеющих коды заказа ЮНИТ-М519-ДЗТ-Р102-К103-К103-К103-В101-В101-В101-В101-В121-М186.0800-М186.0800-С01.03 и ЮНИТ-М514-ДЗО-Р102-К103-К103-К103-х-В101-В101-М1С4.0С00-С01.03}%
  \renewcommand{\functions}{%
    \item дифференциальная защита трансформатора (ДЗТ);
    \item газовая защита автотрансформатора (ГЗ АТ);
    \item технологические защиты автотрансформатора (ТЗ АТ);
	\item газовая защита линейного регулировочного трансформатора (ГЗ ЛРТ);
	\item газовая защита устройство регулирования напряжения под нагрузкой (ГЗ РПН);
    \item контроль изоляции газовых защит (КИ ГЗ);
	\item контроль изоляции технологических защит (КИ ТЗ);
	\item максимальная токовая защита с компбинированным пуском по напряжению стороны НН (МТЗ/U НН);
    \item контроль изоляции стороны НН (КИ НН);	
	\item блокировка при неисправности цепей напряжения стороны НН (БНН НН);	
	\item защита от перегрузки стороны ВН (ЗП ВН);
	\item защита от перегрузки общей обмотки (ЗП ОО);
	\item токовый контроль для защиты от дуговых замыканий (ТК ЗДЗ).}%
  \renewcommand{\ManualUnitDefinite}{ЮТКБ.656122.611-06.02 РЭ «Основная защита автотрансформатора напряжением 220~кВ с функциями УРОВ и АУВ. Устройство типа ЮНИТ-М500-ДЗТ» и ЮТКБ.656122.611-05.01 РЭ «Дифференциальная защита ошиновки от 6 до 750~кВ. Устройство типа ЮНИТ-М500-ДЗО»}%
}

\newcommand{\ManualTwentyFour}{%
  \renewcommand{\funcCabinet}{основных защит автотрансформатора 220~кВ на электрических станциях и подстанциях архитектуры I типа}%
  \renewcommand{\STO}{СТО 56947007-33.040.20.276-2019 ПАО «ФСК ЕЭС»}%
  \renewcommand{\units}{одного интеллектуального электронного устройства (ИЭУ) серии «ЮНИТ-М500», имеющего код заказа ЮНИТ-М519-ДЗТ-Р102-К103-К103-К103-В101-В101-В101-В101-В121-М186.0800-М186.0800-С01.03}%
  \renewcommand{\functions}{%
    \item дифференциальная защита трансформатора (ДЗТ);
    \item газовая защита автотрансформатора (ГЗ АТ);
    \item технологические защиты автотрансформатора (ТЗ АТ);
	\item газовая защита линейного регулировочного трансформатора (ГЗ ЛРТ);
	\item газовая защита устройство регулирования напряжения под нагрузкой (ГЗ РПН);
    \item контроль изоляции газовых защит (КИ ГЗ);
	\item контроль изоляции технологических защит (КИ ТЗ);
	\item максимальная токовая защита с компбинированным пуском по напряжению стороны НН (МТЗ/U НН);
    \item контроль изоляции стороны НН (КИ НН);	
	\item блокировка при неисправности цепей напряжения стороны НН (БНН НН);	
	\item защита от перегрузки стороны ВН (ЗП ВН);
	\item защита от перегрузки общей обмотки (ЗП ОО);
	\item защита от перегрузки стороны НН (ЗП НН);
	\item контроль отсутствия напряжения автотрансформатора (КОН АТ);	
	\item токовый контроль для защиты от дуговых замыканий (ТК ЗДЗ);
	\item устройство резервирования отказа выключателя НН (ПО УРОВ НН);
	\item реле тока пуска охлаждения (РТ ПО).}%
  \renewcommand{\ManualUnitDefinite}{ЮТКБ.656122.611-06.02 РЭ «Основная защита автотрансформатора напряжением 220~кВ с функциями УРОВ и АУВ. Устройство типа ЮНИТ-М500-ДЗТ»}%
}

\newcommand{\ManualTwentyFive}{%
  \renewcommand{\funcCabinet}{резервных защит автотрансформатора 220-750~кВ на электрических станциях и подстанциях архитектуры I типа}%
  \renewcommand{\STO}{СТО 56947007-33.040.20.276-2019 ПАО «ФСК ЕЭС»}%
  \renewcommand{\units}{одного интеллектуального электронного устройства (ИЭУ) серии «ЮНИТ-М500», имеющего код заказа ЮНИТ-М519-РЗТ-Р102-К103-К103-К103-К103-К103-К103-К103-В101-В101-В101-М168.0600-С01.03}%
  \renewcommand{\functions}{%
	\item дистанционная защита (ДЗ);
	\item токовая защита нулевой последовательности (ТЗНП);
	\item максимальная токовая защита (МТЗ);
	\item оперативное ускорение (ОУ);
	\item автоматическое ускорение (АУ);
	\item блокировка при неисправности цепей напряжения (БНН);
	\item блокировка при качаниях (БК);
	\item защита от неполнофазного режима (ЗНР);
	\item блокировка при неисправности цепей напряжения стороны НН (БНН НН).}%
  \renewcommand{\ManualUnitDefinite}{ЮТКБ.656122.611-07.02 РЭ «Резервные защиты автотрансформатора напряжением 220 кВ. Устройство типа ЮНИТ-М500-РЗТ»}%
}

\newcommand{\ManualTwentySix}{%
  \renewcommand{\funcCabinet}{защит батареи статических конденсаторов 110-330~кВ на электрических станциях и подстанциях архитектуры I типа}%
  \renewcommand{\STO}{СТО 56947007-33.040.20.279-2019 ПАО «ФСК ЕЭС»}%
  \renewcommand{\units}{одного интеллектуального электронного устройства (ИЭУ) серии «ЮНИТ-М500», имеющего код заказа ЮНИТ-М514-БСК-Р102-К103-К103-В101-В101-М173.0700-М173.0700-С01.03}%
  \renewcommand{\functions}{%
	\item дифференциальная токовая защита (ДТЗ);
	\item дифференциальная токовая защита по нулевой последовательности (ДТЗ НП);
	\item защита от внутренних повреждений (ЗВП);
	\item защита от перегрузки токами высших гармоник(ЗПВГ);
	\item защита минимального напряжения (ЗМН);
	\item защита от повышения напряжения (ЗПН);
	\item максимальная токовая защита стороны ВН (МТЗ ВН);
	\item токовая защита обратной последовательности стороны ВН (ТЗОП ВН);
	\item токовая защита нулевой последовательности стороны ВН (ТЗНП ВН);
	\item токовая защита нулевой последовательности нейтрали (ТЗНП НЕЙТР);
	\item защита от неполнофазного режима (ЗНР);
	\item блокировка при неисправности цепей напряжения (БНН).}%
  \renewcommand{\ManualUnitDefinite}{ЮТКБ.656122.611-09.01 РЭ «Защиты батареи статических конденсаторов с функциями УРОВ и АУВ. Устройство типа ЮНИТ-М500-БСК»}%
}

\newcommand{\ManualTwentySeven}{%
  \renewcommand{\funcCabinet}{защит батареи статических конденсаторов 110-330~кВ на электрических станциях и подстанциях архитектуры I типа}%
  \renewcommand{\STO}{СТО 56947007-33.040.20.279-2019 ПАО «ФСК ЕЭС»}%
  \renewcommand{\units}{двух интеллектуальных электронных устройств (ИЭУ) серии «ЮНИТ-М500», имеющих коды заказа ЮНИТ-М514-БСК-Р102-К103-К103-В101-В101-М173.0700-М173.0700-С01.03 и ЮНИТ-М514-БСК-Р102-К103-К103-х-В101-В101-В101-М141.0400-С01.03}%
  \renewcommand{\functions}{%
	\item дифференциальная токовая защита (ДТЗ);
	\item дифференциальная токовая защита по нулевой последовательности (ДТЗ НП);
	\item защита от внутренних повреждений (ЗВП);
	\item защита от перегрузки токами высших гармоник(ЗПВГ);
	\item защита минимального напряжения (ЗМН);
	\item защита от повышения напряжения (ЗПН);
	\item максимальная токовая защита стороны ВН (МТЗ ВН);
	\item токовая защита обратной последовательности стороны ВН (ТЗОП ВН);
	\item токовая защита нулевой последовательности стороны ВН (ТЗНП ВН);
	\item токовая защита нулевой последовательности нейтрали (ТЗНП НЕЙТР);
	\item устройство резервирования отказа выключателя (УРОВ);
	\item защита от непереключения фаз (ЗНФ);
    \item автоматика управления выключателем (АУВ).}%
  \renewcommand{\ManualUnitDefinite}{ЮТКБ.656122.611-09.01 РЭ «Защиты батареи статических конденсаторов с функциями УРОВ и АУВ. Устройство типа ЮНИТ-М500-БСК»}%
}

\newcommand{\ManualTwentyEight}{%
  \renewcommand{\funcCabinet}{защит шунтирующего реактора 220-750~кВ на электрических станциях и подстанциях архитектуры I типа}%
  \renewcommand{\STO}{СТО 56947007-33.040.20.279-2019 ПАО «ФСК ЕЭС»}%
  \renewcommand{\units}{одного интеллектуального электронного устройства (ИЭУ) серии «ЮНИТ-М500», имеющего код заказа ЮНИТ-М519-ШР-Р102-К103-К103-К103-В101-В101-В101-В101-В101-М186.0800-М186.0800-С01.03}%
  \renewcommand{\functions}{%
	\item дифференциальная защита реактора (ДЗР);
	\item поперечная дифференциальная защита реактора (ПДЗР);
	\item газовые защиты (ГЗ);
	\item технологические защиты (ТЗ);
	\item контроль изоляции газовых защит (КИ ГЗ);
	\item контроль изоляции технологических защит (КИ ТЗ);
	\item контроль изоляции ввода (КИВ);
	\item максимальная токовая защита стороны ВН (МТЗ ВН);
	\item токовая защита обратной последовательности стороны ВН (ТЗОП ВН);
	\item токовая защита нулевой последовательности стороны ВН (ТЗНП ВН);
	\item токовая защита нулевой последовательности нейтральных выводов(ТЗНП НВ);
	\item защита от неполнофазного режима (ЗНР);
	\item контроль отсутствия напряжения шунтирующего реактора (КОН ШР);
    \item токовое реле (орган) пуска охлаждения (РТПО).}%
  \renewcommand{\ManualUnitDefinite}{ЮТКБ.656122.611-11.01 РЭ «Защиты шунтирующего реактора с функциями УРОВ и АУВ. Устройство типа ЮНИТ-М500-ШР»}%
}

\newcommand{\ManualTwentyNine}{%
  \renewcommand{\funcCabinet}{защит шунтирующего реактора 110-330~кВ на электрических станциях и подстанциях архитектуры I типа}%
  \renewcommand{\STO}{СТО 56947007-33.040.20.279-2019 ПАО «ФСК ЕЭС»}%
  \renewcommand{\units}{одного интеллектуального электронного устройства (ИЭУ) серии «ЮНИТ-М500», имеющего код заказа ЮНИТ-М519-ШР-Р102-К103-К103-К103-В101-В101-В101-В101-В101-М186.0800-М186.0800-С01.03}%
  \renewcommand{\functions}{%
	\item дифференциальная защита реактора (ДЗР);
	\item поперечная дифференциальная защита реактора (ПДЗР);
	\item газовые защиты (ГЗ);
	\item технологические защиты (ТЗ);
	\item контроль изоляции газовых защит (КИ ГЗ);
	\item контроль изоляции технологических защит (КИ ТЗ);
	\item контроль изоляции ввода (КИВ);
	\item максимальная токовая защита стороны ВН (МТЗ ВН);
	\item токовая защита обратной последовательности стороны ВН (ТЗОП ВН);
	\item токовая защита нулевой последовательности стороны ВН (ТЗНП ВН);
	\item токовая защита нулевой последовательности нейтральных выводов(ТЗНП НВ);
	\item защита от неполнофазного режима (ЗНР);
	\item контроль отсутствия напряжения шунтирующего реактора (КОН ШР);
    \item токовое реле (орган) пуска охлаждения (РТПО).}%
  \renewcommand{\ManualUnitDefinite}{ЮТКБ.656122.611-11.01 РЭ «Защиты шунтирующего реактора с функциями УРОВ и АУВ. Устройство типа ЮНИТ-М500-ШР»}%
}

\newcommand{\ManualThirty}{%
  \renewcommand{\funcCabinet}{защит компенсационной обмотки, ошиновки НН и обмотки управления УШРП на электрических станциях и подстанциях напряжением 110-220~кВ архитектуры I типа}%
  \renewcommand{\STO}{СТО 56947007-33.040.20.279-2019 ПАО «ФСК ЕЭС»}%
  \renewcommand{\units}{одного интеллектуального электронного устройства (ИЭУ) серии «ЮНИТ-М500», имеющего код заказа ЮНИТ-М514-ШР-Р102-К103-К103-В101-В101-М186.0800-М186.0800-С01.03}%
  \renewcommand{\functions}{%
	\item максимальная токовая защита компенсационной обмотки реактора (МТЗ КОР);
	\item токовая защита обратной последовательности компенсационной обмотки реактора (ТЗОП КОР);
	\item токовая защита нулевой последовательности стороны компенсационной обмотки реактора (ТЗНП КОР);
	\item максимальная токовая защита стороны НН (МТЗ НН);
	\item контроль изоляции стороны НН (КИ НН);
	\item блокировка при неисправности цепей напряжения стороны НН (БНН НН);
	\item дифференциальная защита ошиновки стороны НН, компенсационной обмотки реактора (ДЗО НН, КОР);
	\item максимальная токовая защита обмотки управления реактора(МТЗ ОУР);
	\item токовый контроль защиты от дуговых замыканий (ТК ЗДЗ).}%
  \renewcommand{\ManualUnitDefinite}{ЮТКБ.656122.611-11.02 РЭ «Защиты обмотки управления, компенсационной обмотки и ошиновки низкой стороны УШРП. Устройство типа ЮНИТ-М500-ШР»}%
}

\newcommand{\ManualThirtyOne}{%
  \renewcommand{\funcCabinet}{дифференциальной защиты для одной или двух секций сборных шин РУ напряжением 35 кВ с жесткой фиксацией присоединений к шинам и количеством не более 10 (включая СВ между защищаемыми секциями) на электрических станциях и подстанциях архитектуры I типа}%
  \renewcommand{\STO}{СТО 34.01-4.1-020-2025 ПАО «ФСК ЕЭС»}%
  \renewcommand{\units}{одного интеллектуального электронного устройства (ИЭУ) серии «ЮНИТ-М500», имеющего код заказа ЮНИТ-М519-ДЗШ-Р102-К103-К103-К103-х-х-В101-В101-В101-В101-В101-М1С4.0С00-С01.03}%
  \renewcommand{\functions}{%
	\item дифференциальная зазита шин (ДЗШ).}%
  \renewcommand{\ManualUnitDefinite}{ЮТКБ.656122.611-04.01 РЭ «Дифференциальная защита сборных шин для первичных схем до 12 элементов на напряжение от 110 до 220 кВ с функцией УРОВ. Устройство типа ЮНИТ-М500-ДЗШ»}%
}

\newcommand{\ManualThirtyTwo}{%
  \renewcommand{\funcCabinet}{дифференциальной защиты для двух секций сборных шин или двух систем сборных шин (не более 9 присоединений; контролируемых групп токов -- 11) на электрических станциях и подстанциях напряжением 110–220~кВ архитектуры II типа}%
  \renewcommand{\STO}{СТО 56947007-33.040.20.286-2019 ПАО «ФСК ЕЭС»}%
  \renewcommand{\units}{одного интеллектуального электронного устройства (ИЭУ) серии «ЮНИТ-М500», имеющего код заказа ЮНИТ-М510-ДЗШ-Р102-К101-х-В101-М1С4.0С00-С01.03}%
  \renewcommand{\functions}{%
	\item дифференциальная зазита шин (ДЗШ);
	\item контроль цепей тока (КЦТ);
    \item блокировка при неисправности цепей напряжения (БНН);
	\item устройство резервирования отказа выключателя (УРОВ).}%
  \renewcommand{\ManualUnitDefinite}{ЮТКБ.656122.611-04.01 РЭ «Дифференциальная защита сборных шин для первичных схем до 12 элементов на напряжение от 110 до 220 кВ с функцией УРОВ. Устройство типа ЮНИТ-М500-ДЗШ»}%
}

\newcommand{\ManualThirtyThree}{%
  \renewcommand{\funcCabinet}{дифференциальной защиты для двух секций сборных шин или двух систем сборных шин (не более 16 присоединений; контролируемых групп токов -- 20) на электрических станциях и подстанциях напряжением 110–220~кВ архитектуры II типа}%
  \renewcommand{\STO}{СТО 56947007-33.040.20.286-2019 ПАО «ФСК ЕЭС»}%
  \renewcommand{\units}{одного интеллектуального электронного устройства (ИЭУ) серии «ЮНИТ-М500», имеющего код заказа ЮНИТ-М514-ДЗШ-Р102-К101-х-В101-В101-М1С4.0С00-М1С4.0С00-С01.03}%
  \renewcommand{\functions}{%
	\item дифференциальная зазита шин (ДЗШ);
	\item контроль цепей тока (КЦТ);
    \item блокировка при неисправности цепей напряжения (БНН);
	\item устройство резервирования отказа выключателя (УРОВ).}%
  \renewcommand{\ManualUnitDefinite}{ЮТКБ.656122.611-04.02 РЭ «Дифференциальная защита сборных шин для первичных схем до 24 элементов на напряжение от 110 до 220 кВ с функцией УРОВ. Устройство типа ЮНИТ-М500-ДЗШ»}%
}

\newcommand{\ManualThirtyFour}{%
  \renewcommand{\funcCabinet}{защиты ошиновок ВН, СН, НН (авто)трансформаторов и ошиновок оборудования СКРМ на подстанциях на электрических станциях и подстанциях напряжением 6–750~кВ архитектуры II типа}%
  \renewcommand{\STO}{СТО 56947007-33.040.20.286-2019 ПАО «ФСК ЕЭС»}%
  \renewcommand{\units}{двух интеллектуальных электронных устройств (ИЭУ) серии «ЮНИТ-М500», имеющих код заказа ЮНИТ-М514-ДЗО-Р102-К101-х-х-В101-М166.0600-М166.0600-С01.03}%
  \renewcommand{\functions}{%
	\item дифференциальная зазита ошиновки (ДЗО);
	\item контроль цепей тока (КЦТ);
    \item блокировка при неисправности цепей напряжения (БНН);
	\item устройство резервирования отказа выключателя В1(В2) (УРОВ В1(В2));
	\item автоматическое повторное включение В1(В2) (АПВ В1(В2));
	\item контроль синхронизма и напряжений В1(В2) (КСН В1(В2));
	\item защита от непереключения фаз В1(В2) (ЗНФ В1(В2));
	\item автоматика управления выключателем В1(В2) (АУВ В1(В2)).}%
  \renewcommand{\ManualUnitDefinite}{ЮТКБ.656122.611-05.01 РЭ «Дифференциальная защита ошиновки от 6 до 750~кВ. Устройство типа ЮНИТ-М500-ДЗО»}%
}

\newcommand{\ManualThirtyFive}{%
  \renewcommand{\funcCabinet}{защиты ошиновок ВН, СН, НН (авто)трансформаторов и ошиновок оборудования СКРМ на подстанциях на электрических станциях и подстанциях напряжением 6–750~кВ архитектуры II типа}%
  \renewcommand{\STO}{СТО 56947007-33.040.20.286-2019 ПАО «ФСК ЕЭС»}%
  \renewcommand{\units}{одного интеллектуального электронного устройства (ИЭУ) серии «ЮНИТ-М500», имеющего код заказа ЮНИТ-М514-ДЗО-Р102-К101-х-х-В101-М166.0600-М166.0600-С01.03}%
  \renewcommand{\functions}{%
	\item дифференциальная зазита ошиновки (ДЗО);
	\item контроль цепей тока (КЦТ);
    \item блокировка при неисправности цепей напряжения (БНН);
	\item устройство резервирования отказа выключателя В1(В2) (УРОВ В1(В2));
	\item автоматическое повторное включение В1(В2) (АПВ В1(В2));
	\item контроль синхронизма и напряжений В1(В2) (КСН В1(В2));
	\item защита от непереключения фаз В1(В2) (ЗНФ В1(В2));
	\item автоматика управления выключателем В1(В2) (АУВ В1(В2)).}%
  \renewcommand{\ManualUnitDefinite}{ЮТКБ.656122.611-05.01 РЭ «Дифференциальная защита ошиновки от 6 до 750~кВ. Устройство типа ЮНИТ-М500-ДЗО»}%
}

\newcommand{\ManualThirtySix}{%
  \renewcommand{\funcCabinet}{ступенчатых защит, УРОВ и АУВ ОВ на электрических станциях и подстанциях напряжением 110–220~кВ архитектуры II типа}%
  \renewcommand{\STO}{СТО 56947007-33.040.20.286-2019 ПАО «ФСК ЕЭС»}%
  \renewcommand{\units}{одного интеллектуального электронного устройства (ИЭУ) серии «ЮНИТ-М500», имеющего код заказа ЮНИТ-М510-ЛВ-Р102-К101-х-В101-М179.0700-С01.03}%
  \renewcommand{\functions}{%
	\item дистанционная защита (ДЗ);
	\item токовая защита нулевой последовательности (ТЗНП);
	\item максимальная токовая защита (МТЗ);
	\item междуфазная токовая отсечка (МФТО);
	\item оперативное ускорение (ОУ);
	\item автоматическое ускорение (АУ);
	\item блокировка при неисправности цепей напряжения (БНН);
	\item блокировка при качаниях (БК);
	\item защита от неполнофазного режима (ЗНР);
	\item определение места повреждения (ОМП);
	\item устройство резервирования отказа выключателя (УРОВ);
	\item автоматическое повторное включение (АПВ);
	\item контроль синхронизма и напряжений (КСН);
	\item защита от непереключения фаз (ЗНФ);
	\item автоматика управления выключателем (АУВ).}%
  \renewcommand{\ManualUnitDefinite}{ЮТКБ.656122.611-03.01 РЭ «Комплект ступенчатых защит для линии напряжением от 110 до 220~кВ с функциями УРОВ и АУВ. Устройство типа ЮНИТ-М500-ЛВ»}%
}	

\newcommand{\ManualThirtySeven}{%
  \renewcommand{\funcCabinet}{ступенчатых защит, УРОВ и АУВ ШСВ/СВ на электрических станциях и подстанциях напряжением 110–750~кВ архитектуры II типа}%
  \renewcommand{\STO}{СТО 56947007-33.040.20.286-2019 ПАО «ФСК ЕЭС»}%
  \renewcommand{\units}{одного интеллектуального электронного устройства (ИЭУ) серии «ЮНИТ-М500», имеющего код заказа ЮНИТ-М510-СВ-Р102-К101-х-В101-М139.0300-С01.03}%
  \renewcommand{\functions}{%
	\item дистанционная защита (ДЗ);
	\item токовая защита нулевой последовательности (ТЗНП);
	\item максимальная токовая защита (МТЗ);
	\item оперативное ускорение (ОУ);
	\item автоматическое ускорение (АУ);
	\item блокировка при неисправности цепей напряжения (БНН);
	\item блокировка при качаниях (БК);
	\item защита от неполнофазного режима (ЗНР);
	\item автоматика опережающего деления сети (АОДС);
	\item определение места повреждения (ОМП);
	\item устройство резервирования отказа выключателя (УРОВ);
	\item автоматическое повторное включение (АПВ);
	\item контроль синхронизма и напряжений (КСН);
	\item защита от непереключения фаз (ЗНФ);
	\item автоматика управления выключателем (АУВ).}%
  \renewcommand{\ManualUnitDefinite}{ЮТКБ.656122.611-08.01 РЭ «Комплект защит шиносоединительного/секционного выключателя с функциями УРОВ и АУВ. Устройство типа ЮНИТ-М500-СВ»}%
}	

\newcommand{\ManualThirtyEight}{%
  \renewcommand{\funcCabinet}{дифференциально-фазной защиты ЛЭП 110-220~кВ с комплектом ступенчатых защит и обменом разрешающими сигналами, УРОВ и АУВ двух выключателей на электрических станциях и подстанциях архитектуры II типа}%
  \renewcommand{\STO}{СТО 56947007-33.040.20.283-2019 ПАО «ФСК ЕЭС»}%
  \renewcommand{\units}{одного интеллектуального электронного устройства (ИЭУ) серии «ЮНИТ-М500», имеющего код заказа ЮНИТ-М510-ВЧ-Р102-К101-В120-М179.0700-В101-С01.03}%
  \renewcommand{\functions}{%
	\item дифференциально-фазная защита (ДФЗ);
	\item дистанционная защита (ДЗ);
	\item токовая защита нулевой последовательности (ТЗНП);
	\item максимальная токовая защита (МТЗ);
	\item междуфазная токовая отсечка (МФТО);
	\item оперативное ускорение (ОУ);
	\item автоматическое ускорение (АУ);
	\item блокировка при неисправности цепей напряжения (БНН);
	\item блокировка при качаниях (БК);
	\item ЛС (осуществляет передачу дискретных сигналов ТУ/ТО по ВЧ каналу с использованием возможностей ВЧ приемопередатчика);
	\item защита от неполнофазного режима (ЗНР);
	\item определение места повреждения (ОМП);
	\item устройство резервирования отказа выключателя В1(В2) (УРОВ В1(В2));
	\item автоматическое повторное включение В1(В2) (АПВ В1(В2));
	\item контроль синхронизма и напряжений В1(В2) (КСН В1(В2));
	\item защита от непереключения фаз В1(В2) (ЗНФ В1(В2));
	\item автоматика управления выключателем В1(В2) (АУВ В1(В2)).}%
  \renewcommand{\ManualUnitDefinite}{ЮТКБ.656122.611-01.01 РЭ «Высокочастотные защиты с комплектом ступенчатых защит для линии напряжением от 110 до 220~кВ с функциями УРОВ и АУВ. Устройство типа ЮНИТ-М500-ВЧ»}%
}

\newcommand{\ManualThirtyNine}{%
  \renewcommand{\funcCabinet}{направленной высокочастотной защиты ЛЭП 110-220~кВ с комплектом ступенчатых защит и обменом разрешающими сигналами, УРОВ и АУВ двух выключателей на электрических станциях и подстанциях архитектуры II типа}%
  \renewcommand{\STO}{СТО 56947007-33.040.20.283-2019 ПАО «ФСК ЕЭС»}%
  \renewcommand{\units}{одного интеллектуального электронного устройства (ИЭУ) серии «ЮНИТ-М500», имеющего код заказа ЮНИТ-М510-ВЧ-Р102-К101-В120-М179.0700-В101-С01.03}%
  \renewcommand{\functions}{%
	\item направлення высокочастотная защита (НВЧЗ);
	\item дистанционная защита (ДЗ);
	\item токовая защита нулевой последовательности (ТЗНП);
	\item максимальная токовая защита (МТЗ);
	\item междуфазная токовая отсечка (МФТО);
	\item оперативное ускорение (ОУ);
	\item автоматическое ускорение (АУ);
	\item блокировка при неисправности цепей напряжения (БНН);
	\item блокировка при качаниях (БК);
	\item ЛС (осуществляет передачу дискретных сигналов ТУ/ТО по ВЧ каналу с использованием возможностей ВЧ приемопередатчика);
	\item защита от неполнофазного режима (ЗНР);
	\item определение места повреждения (ОМП);
	\item устройство резервирования отказа выключателя В1(В2) (УРОВ В1(В2));
	\item автоматическое повторное включение В1(В2) (АПВ В1(В2));
	\item контроль синхронизма и напряжений В1(В2) (КСН В1(В2));
	\item защита от непереключения фаз В1(В2) (ЗНФ В1(В2));
	\item автоматика управления выключателем В1(В2) (АУВ В1(В2)).}%
  \renewcommand{\ManualUnitDefinite}{ЮТКБ.656122.611-01.01 РЭ «Высокочастотные защиты с комплектом ступенчатых защит для линии напряжением от 110 до 220~кВ с функциями УРОВ и АУВ. Устройство типа ЮНИТ-М500-ВЧ»}%
}

\newcommand{\ManualFourty}{%
  \renewcommand{\funcCabinet}{дифференциальной защиты ЛЭП 110-220~кВ с комплектом ступенчатых защит и обменом разрешающими сигналами, УРОВ и АУВ двух выключателей на электрических станциях и подстанциях архитектуры II типа}%
  \renewcommand{\STO}{СТО 56947007-33.040.20.283-2019 ПАО «ФСК ЕЭС»}%
  \renewcommand{\units}{одного интеллектуального электронного устройства (ИЭУ) серии «ЮНИТ-М500», имеющего код заказа ЮНИТ-М510-ДЗЛ-Р102-К101-х-В101-М179.0700-С04.03}%
  \renewcommand{\functions}{%
	\item дифференциальная защита ЛЭП (ДЗЛ);
	\item дистанционная защита (ДЗ);
	\item токовая защита нулевой последовательности (ТЗНП);
	\item максимальная токовая защита (МТЗ);
	\item междуфазная токовая отсечка (МФТО);
	\item оперативное ускорение (ОУ);
	\item автоматическое ускорение (АУ);
	\item блокировка при неисправности цепей напряжения (БНН);
	\item блокировка при качаниях (БК);
	\item ЛС (осуществляет передачу дискретных сигналов ТУ/ТО по ВЧ каналу с использованием возможностей ВЧ приемопередатчика);
	\item защита от неполнофазного режима (ЗНР);
	\item определение места повреждения (ОМП);
	\item устройство резервирования отказа выключателя В1(В2) (УРОВ В1(В2));
	\item автоматическое повторное включение В1(В2) (АПВ В1(В2));
	\item контроль синхронизма и напряжений В1(В2) (КСН В1(В2));
	\item защита от непереключения фаз В1(В2) (ЗНФ В1(В2));
	\item автоматика управления выключателем В1(В2) (АУВ В1(В2)).}%
  \renewcommand{\ManualUnitDefinite}{ЮТКБ.656122.611-02.01 РЭ «Дифференциальная защита линии с комплектом ступенчатых защит для линии напряжением от 110 до 220~кВ с функциями УРОВ и АУВ. Устройство типа ЮНИТ-М500-ДЗЛ»}%
}

\newcommand{\ManualFourtyOne}{%
  \renewcommand{\funcCabinet}{ступенчатых защит ЛЭП 110-220~кВ и обменом разрешающими сигналами, УРОВ и АУВ двух выключателей на электрических станциях и подстанциях архитектуры II типа}%
  \renewcommand{\STO}{СТО 56947007-33.040.20.283-2019 ПАО «ФСК ЕЭС»}%
  \renewcommand{\units}{одного интеллектуального электронного устройства (ИЭУ) серии «ЮНИТ-М500», имеющего код заказа ЮНИТ-М510-ЛВ-Р102-К101-х-В101-М179.0700-С01.03}%
  \renewcommand{\functions}{%
	\item дистанционная защита (ДЗ);
	\item токовая защита нулевой последовательности (ТЗНП);
	\item максимальная токовая защита (МТЗ);
	\item междуфазная токовая отсечка (МФТО);
	\item оперативное ускорение (ОУ);
	\item автоматическое ускорение (АУ);
	\item блокировка при неисправности цепей напряжения (БНН);
	\item блокировка при качаниях (БК);
	\item определение места повреждения (ОМП);
	\item устройство резервирования отказа выключателя В1(В2) (УРОВ В1(В2));
	\item автоматическое повторное включение В1(В2) (АПВ В1(В2));
	\item контроль синхронизма и напряжений В1(В2) (КСН В1(В2));
	\item защита от непереключения фаз В1(В2) (ЗНФ В1(В2));
	\item автоматика управления выключателем В1(В2) (АУВ В1(В2)).}%
  \renewcommand{\ManualUnitDefinite}{ЮТКБ.656122.611-03.01 РЭ «Комплект ступенчатых защит для линии напряжением от 110 до 220~кВ с функциями УРОВ и АУВ. Устройство типа ЮНИТ-М500-ЛВ»}%
}

\newcommand{\ManualFourtyTwo}{%
  \renewcommand{\funcCabinet}{защит трансформатора 110-220~кВ и ошиновки 6-35~кВ, УРОВ, АУВ на электрических станциях и подстанциях архитектуры II типа}%
  \renewcommand{\STO}{СТО 56947007-33.040.20.277-2019 ПАО «ФСК ЕЭС»}%
  \renewcommand{\units}{двух интеллектуальных электронных устройств (ИЭУ) серии «ЮНИТ-М500», имеющих коды заказа ЮНИТ-М514-ДЗТ-Р102-К101-х-х-В101-М186.0800-М186.0800-С01.03 и ЮНИТ-М510-ДЗО-Р102-К101-х-В101-М1С4.0С00-С01.03}%
  \renewcommand{\functions}{%
    \item дифференциальная защита трансформатора (ДЗТ);
    \item газовая защита трансформатора (ГЗ Т);
    \item технологические защиты трансформатора (ТЗ Т);
	\item газовая защита устройство регулирования напряжения под нагрузкой (ГЗ РПН);
	\item максимальная токовая защита с компбинированным пуском по напряжению стороны ВН (МТЗ/U ВН);
	\item токовая защита нулевой последовательности ВН (ТЗНП ВН);
	\item максимальная токовая защита с компбинированным пуском по напряжению стороны НН, СН (МТЗ/U НН, СН);
	\item защита от неполнофазного режима (ЗНР);
	\item токовая защита нулевой последовательности резистора заземления нейтрали (ТЗНП РЗН);
    \item контроль изоляции газовых защит (КИ ГЗ);
	\item контроль изоляции технологических защит (КИ ТЗ);
	\item контроль изоляции стороны НН, СН (КИ НН, СН);
	\item блокировка при неисправности цепей напряжения стороны НН, СН (БНН НН, СН);
	\item защита от перегрузки стороны ВН (ЗП ВН);
	\item защита от перегрузки стороны НН, СН (ЗП НН, СН);
	\item контроль отсутствия напряжения трансформатора (КОН Т);
	\item токовый контроль для защиты от дуговых замыканий (ТК ЗДЗ);
	\item устройство резервирования отказа выключателя стороны НН (ПО УРОВ НН);
	\item токовое реле (орган) пуска охлаждения (РТПО);
	\item устройство резервирования отказа выключателя стороны (УРОВ);
	\item автоматическое повторное включение (АПВ);
	\item контроль синхронизма и напряжений (КСН);
	\item автоматика управления выключателем (АУВ);
	\item дифференциальная защита ошиновки стороны НН, СН (ДЗО НН, СН).}%
  \renewcommand{\ManualUnitDefinite}{ЮТКБ.656122.611-06.01 РЭ «Основная защита трансформатора напряжением от 110 до 220~кВ с функциями УРОВ и АУВ. Устройство типа ЮНИТ-М500-ДЗТ» и ЮТКБ.656122.611-05.01 РЭ «Дифференциальная защита ошиновки от 6 до 750~кВ. Устройство типа ЮНИТ-М500-ДЗО»}%
}

\newcommand{\ManualFourtyThree}{%
  \renewcommand{\funcCabinet}{защит трансформатора 110-220~кВ и ошиновки 6-35~кВ, УРОВ, АУВ на электрических станциях и подстанциях архитектуры II типа}%
  \renewcommand{\STO}{СТО 56947007-33.040.20.277-2019 ПАО «ФСК ЕЭС»}%
  \renewcommand{\units}{одного интеллектуального электронного устройства (ИЭУ) серии «ЮНИТ-М500», имеющего код заказа ЮНИТ-М514-ДЗТ-Р102-К101-х-х-В101-М186.0800-М186.0800-С01.03}%
  \renewcommand{\functions}{%
    \item дифференциальная защита трансформатора (ДЗТ);
    \item газовая защита трансформатора (ГЗ Т);
    \item технологические защиты трансформатора (ТЗ Т);
	\item газовая защита устройство регулирования напряжения под нагрузкой (ГЗ РПН);
	\item максимальная токовая защита с компбинированным пуском по напряжению стороны ВН (МТЗ/U ВН);
	\item токовая защита нулевой последовательности ВН (ТЗНП ВН);
	\item максимальная токовая защита с компбинированным пуском по напряжению стороны НН, СН (МТЗ/U НН, СН);
	\item защита от неполнофазного режима (ЗНР);
	\item токовая защита нулевой последовательности резистора заземления нейтрали (ТЗНП РЗН);
	\item блокировка при неисправности цепей напряжения стороны НН, СН (БНН НН, СН);
	\item защита от перегрузки стороны ВН (ЗП ВН);
	\item защита от перегрузки стороны НН, СН (ЗП НН, СН);
	\item контроль отсутствия напряжения трансформатора (КОН Т);
	\item токовый контроль для защиты от дуговых замыканий (ТК ЗДЗ);
	\item устройство резервирования отказа выключателя стороны НН (ПО УРОВ НН);
	\item токовое реле (орган) пуска охлаждения (РТПО);
	\item устройство резервирования отказа выключателя стороны НН (УРОВ);
	\item автоматическое повторное включение (АПВ);
	\item контроль синхронизма и напряжений (КСН);
	\item автоматика управления выключателем (АУВ).}%
  \renewcommand{\ManualUnitDefinite}{ЮТКБ.656122.611-06.01 РЭ «Основная защита трансформатора напряжением от 110 до 220~кВ с функциями УРОВ и АУВ. Устройство типа ЮНИТ-М500-ДЗТ»}%
}

\newcommand{\ManualFourtyFour}{%
  \renewcommand{\funcCabinet}{резервных защит трансформатора 110-220~кВ на электрических станциях и подстанциях архитектуры II типа}%
  \renewcommand{\STO}{СТО 56947007-33.040.20.277-2019 ПАО «ФСК ЕЭС»}%
  \renewcommand{\units}{одного интеллектуального электронного устройства (ИЭУ) серии «ЮНИТ-М500», имеющего код заказа ЮНИТ-М510-РЗТ-Р102-К101-х-В101-М168.0600-С01.03}%
  \renewcommand{\functions}{%
    \item газовая защита трансформатора (ГЗ Т);
    \item технологические защиты трансформатора (ТЗ Т);
	\item газовая защита устройство регулирования напряжения под нагрузкой (ГЗ РПН);
	\item максимальная токовая защита с компбинированным пуском по напряжению стороны ВН (МТЗ/U ВН);
	\item токовая защита нулевой последовательности ВН (ТЗНП ВН);
	\item максимальная токовая защита с компбинированным пуском по напряжению стороны СН (МТЗ/U СН);
	\item защита от неполнофазного режима (ЗНР);
	\item блокировка при неисправности цепей напряжения стороны НН, СН (БНН НН, СН).}%
  \renewcommand{\ManualUnitDefinite}{ЮТКБ.656122.611-07.01 РЭ «Резервные защиты трансформатора напряжением от 110 до 220 кВ. Устройство типа ЮНИТ-М500-РЗТ»}%
}

\newcommand{\ManualFourtyFive}{%
  \renewcommand{\funcCabinet}{основных защит автотрансформатора 220~кВ и ошиновки 6-35~кВ, УРОВ, АУВ на электрических станциях и подстанциях архитектуры II типа}%
  \renewcommand{\STO}{СТО 56947007-33.040.20.277-2019 ПАО «ФСК ЕЭС»}%
  \renewcommand{\units}{двух интеллектуальных электронных устройств (ИЭУ) серии «ЮНИТ-М500», имеющих коды заказа ЮНИТ-М514-ДЗТ-Р102-К101-х-х-В101-М186.0800-М186.0800-С01.03 и ЮНИТ-М510-ДЗО-Р102-К101-х-В101-М1С4.0С00-С01.03}%
  \renewcommand{\functions}{%
    \item дифференциальная защита автотрансформатора (ДЗТ);
    \item газовая защита автотрансформатора (ГЗ АТ);
    \item технологические защиты автотрансформатора (ТЗ АТ);
	\item газовая защита линейного регулировочного трансформатора (ГЗ ЛРТ);
	\item технологические защиты линейного регулировочного трансформатора (ТЗ ЛРТ);
	\item газовая защита устройство регулирования напряжения под нагрузкой (ГЗ РПН);
	\item максимальная токовая защита с компбинированным пуском по напряжению стороны НН (МТЗ/U НН);
	\item контроль изоляции стороны НН (КИ НН);
	\item блокировка при неисправности цепей напряжения стороны НН (БНН НН);
	\item защита от перегрузки стороны ВН (ЗП ВН);
	\item защита от перегрузки общей обмотки (ЗП ОО);
	\item защита от перегрузки стороны НН (ЗП НН);
	\item контроль отсутствия напряжения автотрансформатора (КОН АТ);	
	\item токовый контроль для защиты от дуговых замыканий (ТК ЗДЗ);
	\item устройство резервирования отказа выключателя НН (ПО УРОВ НН);
	\item реле тока пуска охлаждения (РТ ПО);
	\item устройство резервирования отказа выключателя стороны ВН, СН (УРОВ ВН, СН);
	\item автоматическое повторное включение стороны ВН, СН (АПВ ВН, СН);
	\item контроль синхронизма и напряжений стороны ВН, СН (КСН ВН, СН);
	\item автоматика управления выключателем стороны ВН, СН (АУВ ВН, СН);
	\item дифференциальная защита ошиновки стороны НН, СН (ДЗО НН, СН).}%
  \renewcommand{\ManualUnitDefinite}{ЮТКБ.656122.611-06.02 РЭ «Основная защита автотрансформатора напряжением 220~кВ с функциями УРОВ и АУВ. Устройство типа ЮНИТ-М500-ДЗТ» и ЮТКБ.656122.611-05.01 РЭ «Дифференциальная защита ошиновки от 6 до 750~кВ. Устройство типа ЮНИТ-М500-ДЗО»}%
}

\newcommand{\ManualFourtySix}{%
  \renewcommand{\funcCabinet}{основных защит автотрансформатора 220~кВ, УРОВ, АУВ на электрических станциях и подстанциях архитектуры II типа}%
  \renewcommand{\STO}{СТО 56947007-33.040.20.277-2019 ПАО «ФСК ЕЭС»}%
  \renewcommand{\units}{одного интеллектуального электронного устройства (ИЭУ) серии «ЮНИТ-М500», имеющего код заказа ЮНИТ-М514-ДЗТ-Р102-К101-х-х-В101-М186.0800-М186.0800-С01.03}%
  \renewcommand{\functions}{%
    \item дифференциальная защита автотрансформатора (ДЗТ);
    \item газовая защита автотрансформатора (ГЗ АТ);
    \item технологические защиты автотрансформатора (ТЗ АТ);
	\item газовая защита линейного регулировочного трансформатора (ГЗ ЛРТ);
	\item технологические защиты линейного регулировочного трансформатора (ТЗ ЛРТ);
	\item газовая защита устройство регулирования напряжения под нагрузкой (ГЗ РПН);
	\item максимальная токовая защита с компбинированным пуском по напряжению стороны НН (МТЗ/U НН);
	\item контроль изоляции стороны НН (КИ НН);
	\item защита от перегрузки стороны ВН (ЗП ВН);
	\item защита от перегрузки общей обмотки (ЗП ОО);
	\item защита от перегрузки стороны НН (ЗП НН);
	\item контроль отсутствия напряжения автотрансформатора (КОН АТ);	
	\item токовый контроль для защиты от дуговых замыканий (ТК ЗДЗ);
	\item устройство резервирования отказа выключателя НН (ПО УРОВ НН);
	\item реле тока пуска охлаждения (РТ ПО);
	\item устройство резервирования отказа выключателя стороны ВН, СН (УРОВ ВН, СН);
	\item автоматическое повторное включение стороны ВН, СН (АПВ ВН, СН);
	\item контроль синхронизма и напряжений стороны ВН, СН (КСН ВН, СН);
	\item автоматика управления выключателем стороны ВН, СН (АУВ ВН, СН);
	\item дифференциальная защита ошиновки стороны НН, СН (ДЗО НН, СН).}%
  \renewcommand{\ManualUnitDefinite}{ЮТКБ.656122.611-06.02 РЭ «Основная защита автотрансформатора напряжением 220~кВ с функциями УРОВ и АУВ. Устройство типа ЮНИТ-М500-ДЗТ»}%
}

\newcommand{\ManualFourtySeven}{%
  \renewcommand{\funcCabinet}{резервных защит автотрансформатора 220-750~кВ на электрических станциях и подстанциях архитектуры II типа}%
  \renewcommand{\STO}{СТО 56947007-33.040.20.277-2019 ПАО «ФСК ЕЭС»}%
  \renewcommand{\units}{двух интеллектуальных электронных устройств (ИЭУ) серии «ЮНИТ-М500», имеющих код заказа ЮНИТ-М510-РЗТ-Р102-К101-х-В101-М168.0600-С01.03}%
  \renewcommand{\functions}{%
    \item дистанционная защита стороны ВН, СН (ДЗ ВН, СН);
	\item токовая защита нулевой последовательности стороны ВН, СН (ТЗНП ВН, СН);
	\item максимальная токовая защита стороны ВН, СН (МТЗ ВН, СН);
	\item оперативное ускорение стороны ВН, СН (ОУ ВН, СН);
	\item автоматическое ускорение стороны ВН, СН (АУ ВН, СН);
	\item блокировка при неисправности цепей напряжения стороны ВН, СН, НН (БНН ВН, СН, НН);
	\item блокировка при качаниях стороны ВН, СН (БК ВН, СН);
	\item защита от неполнофазного режима стороны ВН, СН (ЗНР ВН, СН).}%
  \renewcommand{\ManualUnitDefinite}{ЮТКБ.656122.611-07.02 РЭ «Резервные защиты автотрансформатора напряжением 220 кВ. Устройство типа ЮНИТ-М500-РЗТ»}%
}

\newcommand{\ManualFourtyEight}{%
  \renewcommand{\funcCabinet}{защит батареи статических конденсаторов 110-330~кВ на электрических станциях и подстанциях архитектуры II типа}%
  \renewcommand{\STO}{СТО 56947007-33.040.20.280-2019 ПАО «ФСК ЕЭС»}%
  \renewcommand{\units}{одного интеллектуального электронного устройства (ИЭУ) серии «ЮНИТ-М500», имеющего код заказа ЮНИТ-М514-БСК-Р102-К101-х-х-В101-М173.0700-М173.0700-С01.03}%
  \renewcommand{\functions}{%
	\item дифференциальная токовая защита (ДТЗ);
	\item дифференциальная токовая защита по нулевой последовательности (ДТЗ НП);
	\item защита от внутренних повреждений (ЗВП);
	\item защита от перегрузки токами высших гармоник(ЗПВГ);
	\item защита минимального напряжения (ЗМН);
	\item защита от повышения напряжения (ЗПН);
	\item максимальная токовая защита стороны ВН (МТЗ ВН);
	\item токовая защита обратной последовательности стороны ВН (ТЗОП ВН);
	\item токовая защита нулевой последовательности стороны ВН (ТЗНП ВН);
	\item токовая защита нулевой последовательности нейтрали (ТЗНП НЕЙТР);
	\item защита от неполнофазного режима (ЗНР);
	\item блокировка при неисправности цепей напряжения (БНН);
	\item устройство резервирования отказа выключателя (УРОВ);
	\item защита от непереключения фаз (ЗНФ);
    \item автоматика управления выключателем (АУВ).}%
  \renewcommand{\ManualUnitDefinite}{ЮТКБ.656122.611-09.01 РЭ «Защиты батареи статических конденсаторов с функциями УРОВ и АУВ. Устройство типа ЮНИТ-М500-БСК»}%
}

\newcommand{\ManualFourtyNine}{%
  \renewcommand{\funcCabinet}{защит шунтирующего реактора 110-750~кВ на электрических станциях и подстанциях архитектуры II типа}%
  \renewcommand{\STO}{СТО 56947007-33.040.20.280-2019 ПАО «ФСК ЕЭС»}%
  \renewcommand{\units}{одного интеллектуального электронного устройства (ИЭУ) серии «ЮНИТ-М500», имеющего код заказа ЮНИТ-М514-ШР-Р102-К101-х-х-В101-М186.0800-М186.0800-С01.03}%
  \renewcommand{\functions}{%
	\item дифференциальная защита реактора (ДЗР);
	\item поперечная дифференциальная защита реактора (ПДЗР);
	\item газовые защиты (ГЗ);
	\item технологические защиты (ТЗ);
	\item контроль изоляции газовых защит (КИ ГЗ);
	\item контроль изоляции технологических защит (КИ ТЗ);
	\item контроль изоляции ввода (КИВ);
	\item максимальная токовая защита стороны ВН (МТЗ ВН);
	\item токовая защита обратной последовательности стороны ВН (ТЗОП ВН);
	\item токовая защита нулевой последовательности стороны ВН (ТЗНП ВН);
	\item токовая защита нулевой последовательности нейтральных выводов(ТЗНП НВ);
	\item защита от неполнофазного режима (ЗНР);
	\item контроль отсутствия напряжения шунтирующего реактора (КОН ШР);
    \item токовое реле (орган) пуска охлаждения (РТПО);
	\item автоматика управления реактором (АУР);
	\item устройство резервирования отказа выключателя (УРОВ);
	\item защита от непереключения фаз (ЗНФ);
	\item автоматика управления выключателем (АУВ).}%
  \renewcommand{\ManualUnitDefinite}{ЮТКБ.656122.611-11.01 РЭ «Защиты шунтирующего реактора с функциями УРОВ и АУВ. Устройство типа ЮНИТ-М500-ШР»}%
}

\newcommand{\ManualFifty}{%
  \renewcommand{\funcCabinet}{защит компенсационной обмотки и ее ошиновки, обмотки управления управляемого шунтирующего реактора 110-750~кВ с подмагничиванием на электрических станциях и подстанциях архитектуры II типа}%
  \renewcommand{\STO}{СТО 56947007-33.040.20.280-2019 ПАО «ФСК ЕЭС»}%
  \renewcommand{\units}{одного интеллектуального электронного устройства (ИЭУ) серии «ЮНИТ-М500», имеющего код заказа ЮНИТ-М514-ШР-Р102-К101-х-х-В101-М186.0800-М186.0800-С01.03}%
  \renewcommand{\functions}{%
	\item максимальная токовая защита компенсационной обмотки реактора (МТЗ КОР);
	\item токовая защита обратной последовательности компенсационной обмотки реактора (ТЗОП КОР);
	\item токовая защита нулевой последовательности стороны компенсационной обмотки реактора (ТЗНП КОР);
	\item максимальная токовая защита стороны НН (МТЗ НН);
	\item контроль изоляции стороны НН (КИ НН);
	\item блокировка при неисправности цепей напряжения стороны НН (БНН НН);
	\item дифференциальная защита ошиновки стороны НН, компенсационной обмотки реактора (ДЗО НН, КОР);
	\item максимальная токовая защита обмотки управления реактора(МТЗ ОУР);
	\item токовый контроль защиты от дуговых замыканий (ТК ЗДЗ).}%
  \renewcommand{\ManualUnitDefinite}{ЮТКБ.656122.611-11.02 РЭ «Защиты обмотки управления, компенсационной обмотки и ошиновки низкой стороны УШРП. Устройство типа ЮНИТ-М500-ШР»}%
}

\newcommand{\ManualFiftyOne}{%
  \renewcommand{\funcCabinet}{дифференциальной защиты для одной или двух секций сборных шин 35 кВ с жесткой фиксацией присоединений к шинам и количеством не более 10 (включая СВ между защищаемыми секциями) на электрических станциях и подстанциях архитектуры II типа}%
  \renewcommand{\STO}{«Типовые шкафы ШЭТ РЗА сборных шин, ошиновок и шинных аппаратов 6–750 кВ. Архитектура II типа»}%
  \renewcommand{\units}{одного интеллектуального электронного устройства (ИЭУ) серии «ЮНИТ-М500», имеющего код заказа ЮНИТ-М514-ДЗШ-Р102-К103-К103-х-х-В101-В101-М1С4.0С00-С01.03}%
  \renewcommand{\functions}{%
	\item дифференциальная защита шин (ДЗШ);
	\item устройство резервирования отказа выключателя (УРОВ).}%
  \renewcommand{\ManualUnitDefinite}{ЮТКБ.656122.611-04.01 РЭ «Дифференциальная защита сборных шин для первичных схем до 12 элементов на напряжение от 110 до 220 кВ с функцией УРОВ. Устройство типа ЮНИТ-М500-ДЗШ»}%
}


% Значения по умолчанию (если код не распознан)
\newcommand{\valuesDefault}{%
  \renewcommand{\funcCabinet}{не задано}%
  \renewcommand{\STO}{не задано}%
  \renewcommand{\units}{не задано}%
  \renewcommand{\functions}{\item Не задано}%
}

% =============================================
% Главный макрос выбора в зависимости от \devicecode
% =============================================
\newcommand{\setDeviceVariables}{%
  \IfStrEqCase{\devicecode}{%
    {ЮТКБ.656463.300-001 РЭ}{\ManualOne}		% ШЭТ 240.01-0
	{ЮТКБ.656463.300-002 РЭ}{\ManualTwo}		% ШЭТ 240.02-0
	{ЮТКБ.656463.300-003 РЭ}{\ManualThree}		% ШЭТ 240.03-0
	{ЮТКБ.656463.300-004 РЭ}{\ManualFour}		% ШЭТ 240.04-0
	{ЮТКБ.656463.300-005 РЭ}{\ManualFive}		% ШЭТ 440.01-0
	{ЮТКБ.656463.300-006 РЭ}{\ManualSix}		% ШЭТ 440.02-0
	{ЮТКБ.656463.300-007 РЭ}{\ManualSeven}		% ШЭТ 251.02-0
	{ЮТКБ.656463.300-008 РЭ}{\ManualEight}		% ШЭТ 451.01-0
	{ЮТКБ.656463.300-009 РЭ}{\ManualNine}		% ШЭТ 251.01-0
	{ЮТКБ.656463.300-010 РЭ}{\ManualTen}		% ШЭТ 451.02-0
	{ЮТКБ.656463.300-011 РЭ}{\ManualEleven}		% ШЭТ 220.01-0
	{ЮТКБ.656463.300-012 РЭ}{\ManualTwelve}		% ШЭТ 220.02-0
	{ЮТКБ.656463.300-013 РЭ}{\ManualThirteen}	% ШЭТ 220.03-0
	{ЮТКБ.656463.300-014 РЭ}{\ManualFourteen}	% ШЭТ 220.04-0
	{ЮТКБ.656463.300-015 РЭ}{\ManualFiveteen}	% ШЭТ 220.05-0
	{ЮТКБ.656463.300-016 РЭ}{\ManualSixteen}	% ШЭТ 220.06-0
	{ЮТКБ.656463.300-017 РЭ}{\ManualSeventeen}	% ШЭТ 220.07-0
	{ЮТКБ.656463.300-018 РЭ}{\ManualEighteen}	% ШЭТ 220.08-0
	{ЮТКБ.656463.300-019 РЭ}{\ManualNineteen}	% ШЭТ 221.09-0
    {ЮТКБ.656463.300-020 РЭ}{\ManualTwenty}		% ШЭТ 210.01-0
	{ЮТКБ.656463.300-021 РЭ}{\ManualTwentyOne}	% ШЭТ 210.02-0
	{ЮТКБ.656463.300-022 РЭ}{\ManualTwentyTwo}	% ШЭТ 210.03-0
	{ЮТКБ.656463.300-023 РЭ}{\ManualTwentyThree}% ШЭТ 210.04-0
	{ЮТКБ.656463.300-024 РЭ}{\ManualTwentyFour}	% ШЭТ 210.05-0
	{ЮТКБ.656463.300-025 РЭ}{\ManualTwentyFive}	% ШЭТ 410.01-0
	{ЮТКБ.656463.300-026 РЭ}{\ManualTwentySix}	% ШЭТ 430.01-0
	{ЮТКБ.656463.300-027 РЭ}{\ManualTwentySeven}% ШЭТ 431.02-0
	{ЮТКБ.656463.300-028 РЭ}{\ManualTwentyEight}% ШЭТ 430.03-0
	{ЮТКБ.656463.300-029 РЭ}{\ManualTwentyNine}	% ШЭТ 430.04-0
	{ЮТКБ.656463.300-030 РЭ}{\ManualThirty}		% ШЭТ 430.05-0
	{ЮТКБ.656463.300-031 РЭ}{\ManualThirtyOne}	% ШЭТ 140.01-0
	{ЮТКБ.656463.300-032 РЭ}{\ManualThirtyTwo}	% ШЭТ 240.01-1
	{ЮТКБ.656463.300-033 РЭ}{\ManualThirtyThree}% ШЭТ 240.02-1
	{ЮТКБ.656463.300-034 РЭ}{\ManualThirtyFour}	% ШЭТ 441.01-1
	{ЮТКБ.656463.300-035 РЭ}{\ManualThirtyFive}	% ШЭТ 441.02-1
	{ЮТКБ.656463.300-036 РЭ}{\ManualThirtySix}	% ШЭТ 251.01-1
	{ЮТКБ.656463.300-037 РЭ}{\ManualThirtySeven}% ШЭТ 451.01-1
	{ЮТКБ.656463.300-038 РЭ}{\ManualThirtyEight}% ШЭТ 221.01-1
	{ЮТКБ.656463.300-039 РЭ}{\ManualThirtyNine}	% ШЭТ 221.02-1
	{ЮТКБ.656463.300-040 РЭ}{\ManualFourty}		% ШЭТ 221.03-1
	{ЮТКБ.656463.300-041 РЭ}{\ManualFourtyOne}	% ШЭТ 221.04-1
	{ЮТКБ.656463.300-042 РЭ}{\ManualFourtyTwo}	% ШЭТ 211.01-1
	{ЮТКБ.656463.300-043 РЭ}{\ManualFourtyThree}% ШЭТ 211.02-1
	{ЮТКБ.656463.300-044 РЭ}{\ManualFourtyFour}	% ШЭТ 210.03-1
	{ЮТКБ.656463.300-045 РЭ}{\ManualFourtyFive}	% ШЭТ 211.04-1
	{ЮТКБ.656463.300-046 РЭ}{\ManualFourtySix}	% ШЭТ 211.05-1
	{ЮТКБ.656463.300-047 РЭ}{\ManualFourtySeven}% ШЭТ 410.01-1
	{ЮТКБ.656463.300-048 РЭ}{\ManualFourtyEight}% ШЭТ 431.01-1
	{ЮТКБ.656463.300-049 РЭ}{\ManualFourtyNine}	% ШЭТ 431.02-1
	{ЮТКБ.656463.300-050 РЭ}{\ManualFifty}		% ШЭТ 430.03-1
	{ЮТКБ.656463.300-051 РЭ}{\ManualFiftyOne}	% ШЭТ 140.11-1	
	
    % добавьте другие коды аналогично
  }[\valuesDefault]%
}

\setDeviceVariables


% ============================================================
% РАЗДЕЛ «ОПИСАНИЕ И РАБОТА»
% ============================================================
{\color{unidarkgreen}\section[Описание и работа]{Описание и работа}}

% ============================================================
% 1. НАЗНАЧЕНИЕ
% ============================================================
{\color{unidarkgreen}\subsection{Назначение}}

\begin{enumsub}
\item Шкаф «\devicetype» предназначен для выполнения функций \funcCabinet~при новом строительстве или реконструкции энергообъектов.
\item Шкаф предусмотрен к установке и эксплуатации в помещениях релейных щитов и общеподстанционных пунктов управления (ОПУ).
\item Функциональное назначение шкафа определено структурой его условного обозначения в соответствии с \STO.
\item Исполнение шкафа выполнено с установкой в нем \units~(см. пункт \ref{sec1:4}). \textls[-5]{Индивидуальная карта заказа шкафа представлена в приложении \ref{app:card}.}
\item Шкаф предназначен для выполнения следующих функций релейной защиты и автоматики:
  \begin{itemize}
	\functions
  \end{itemize}
\item Дополнительно шкаф выполняет функции:
  \begin{itemize}
  \item местной предупредительной и аварийной сигнализации;
  \item регистрации аварийных событий (РАС);
  \item самодиагностики.
  \end{itemize}
\end{enumsub}

% ============================================================
% 2. ТЕХНИЧЕСКИЕ ХАРАКТЕРИСТИКИ
% ============================================================
{\color{unidarkgreen}\subsection{Технические характеристики}}

\subsubsection{Основные параметры}

\begin{enumsubsub}
\item Номинальные параметры шкафа указаны в таблице \ref{sec1:tbl1}.
  \begingroup
  \small
  \setlength{\extrarowheight}{0.06cm}
  \begin{longtable}{|>{\centering\arraybackslash}m{12.9cm}|>{\centering\arraybackslash}m{4cm}|}
    \caption{Основные номинальные параметры шкафа\hfill\vspace{-0.5\baselineskip}}\label{sec1:tbl1}\\
    \hline
    \rowcolor{gray!30} Параметр & Значение \\ \hline
    \endfirsthead
    \caption*{\hspace{3pt}\emph{Продолжение таблицы \ref{sec1:tbl1}\hfill\vspace{-0.5\baselineskip}}}\\
    \hline
    \rowcolor{gray!30} Параметр & Значение \\ \hline
    \endhead
    \hline
    \endfoot
    \endlastfoot
    \raggedright Номинальный переменный ток, А & \centering\arraybackslash 1 или 5 \\ \hline
    \raggedright Номинальное линейное напряжение переменного тока, В & \centering\arraybackslash 100 \\ \hline
    \raggedright Номинальная частота, Гц & \centering\arraybackslash 50 \\ \hline
    \raggedright Номинальное напряжение оперативного постоянного тока, В & \centering\arraybackslash 220 \\ \hline
  \end{longtable}
  \endgroup



\IfStrEqCase{\devicecode}{%
  {ЮТКБ.656463.300-020 РЭ}{\def\powerValue{180}}%
  {ЮТКБ.656463.300-023 РЭ}{\def\powerValue{180}}%
  {ЮТКБ.656463.300-042 РЭ}{\def\powerValue{180}}%
  {ЮТКБ.656463.300-045 РЭ}{\def\powerValue{180}}%
  {ЮТКБ.656463.300-005 РЭ}{\def\powerValue{180}}%
  {ЮТКБ.656463.300-034 РЭ}{\def\powerValue{180}}%
  {ЮТКБ.656463.300-047 РЭ}{\def\powerValue{180}}%
  {ЮТКБ.656463.300-027 РЭ}{\def\powerValue{180}}%
}[\def\powerValue{85}]


\item Потребляемая мощность шкафа указана в таблице \ref{sec1:tbl2}.
  \begingroup
  \small
  \setlength{\extrarowheight}{0.06cm}
  \begin{longtable}{|>{\centering\arraybackslash}m{7cm}|>{\centering\arraybackslash}m{5.5cm}|>{\centering\arraybackslash}m{4cm}|}
    \caption{Потребляемая мощность шкафа при подведении к нему номинальных величин\hfill\vspace{-0.5\baselineskip}}\label{sec1:tbl2}\\
    \hline
    \rowcolor{gray!30} \multicolumn{2}{|c|}{Параметр} & Значение (не более) \\ \hline
    \endfirsthead
    \caption*{\hspace{3pt}\emph{Продолжение таблицы \ref{sec1:tbl2}\hfill\vspace{-0.5\baselineskip}}}\\
    \hline
    \rowcolor{gray!30} \multicolumn{2}{|c|}{Параметр} & Значение (не более) \\ \hline
    \endhead
    \hline
    \endfoot
    \endlastfoot
    \multicolumn{2}{|>{\raggedright\arraybackslash}p{12.5cm}|}{По каждому аналоговому каналу при номинальном напряжении, В$\cdot$А} & 0,5 \\ \hline
    \multirow{2}{=}{По каждому аналоговому каналу при номинальном токе, В$\cdot$А} & $I_{\text{ном}} = 1$ А & 0,2 \\ \cline{2-3}
    & $I_{\text{ном}} = 5$ А & 0,5 \\ \hline
    \multicolumn{2}{|>{\raggedright\arraybackslash}p{12.5cm}|}{По цепям напряжения оперативного постоянного тока, Вт} & \powerValue \\ \hline
    \multicolumn{2}{|>{\raggedright\arraybackslash}p{12.5cm}|}{По цепям сигнализации в режиме срабатывания, Вт} & 26 \\ \hline
    \multicolumn{2}{|>{\raggedright\arraybackslash}p{12.5cm}|}{По цепям освещения, Вт} & 7 \\ \hline
  \end{longtable}
  \endgroup
\end{enumsubsub}

\subsubsection{Условия эксплуатации}

\begin{enumsubsub}
\item Климатическое исполнение шкафа и категория его размещения по ГОСТ 15150-69 – УХЛ3.1(УХЛ4).
\item Условия эксплуатации в соответствии с ГОСТ IEC 61439-1-2024:
  \begin{itemize}
  \item температура окружающей среды должна быть не более 40~°С, а средняя температура за 24~ч -- не более 35~°С;
  \item минимальное значение температуры окружающей среды -- минус 5~°С;
  \item \textls[-5]{относительная влажность воздуха не должна превышать 50~\% при максимальной температуре 40~°С. При более низких температурах допускается более высокая относительная влажность, например 90~\% при 20~°С;}
  \item высота установки над уровнем моря не должна превышать 2000~м;
  \item степень загрязнения окружающей среды 1 -- загрязнение отсутствует или имеется только сухое непроводящее загрязнение. Загрязнение незначительное;
  \item место установки шкафа должно быть защищено от попадания брызг воды, масел, эмульсий, а также от прямого воздействия солнечной радиации.
  \end{itemize}
\item Рабочее положение шкафа в пространстве -- вертикальное с отклонением от рабочего положения до 5° в любую сторону.
\item \textls[-20]{В части воздействия факторов внешней среды шкаф удовлетворяет требованиям группы механического исполнения М40 по ГОСТ 30631-99. Уровень вибрационных нагрузок от 0,5~Гц до 100~Гц с ускорением 0,5~g. Шкаф выдерживает многократные ударные нагрузки длительностью от 2~мс до 20~мс с максимальным ускорением 3~g.}
\item Шкаф сейсмостоек при воздействии землетрясений интенсивностью до 9 баллов включительно по шкале MSK-64 при уровне установки над нулевой отметкой до 10~м по ГОСТ 17516.1-90.
\item Степень защиты, обеспечиваемая оболочкой шкафа, по ГОСТ 14254-2015 от доступа к опасным частям, попадания внешних твердых предметов и (или) воды соответствует IP54.
\end{enumsubsub}

\subsubsection{Сопротивление и электрическая прочность изоляции}

\begin{enumsubsub}
\item Сопротивление изоляции электрически независимых цепей шкафа, измеренное в холодном состоянии при температуре окружающего воздуха 20±5~°С и относительной влажности до 80~\%, должно быть не менее 100~МОм при напряжении постоянного тока 500~В.
\item В состоянии поставки электрическая изоляция между всеми независимыми цепями шкафа (кроме портов последовательной передачи данных ИЭУ) относительно корпуса и всех независимых цепей между собой выдерживает без пробоя и перекрытия испытательное напряжение, приведенное в таблице \ref{sec1:tbl3}.
  \begingroup
  \small
  \setlength{\extrarowheight}{0.06cm}
  \begin{longtable}{|>{\centering\arraybackslash}m{12.9cm}|>{\centering\arraybackslash}m{4cm}|}
    \caption{Требования к испытанию изоляции\hfill\vspace{-0.5\baselineskip}}\label{sec1:tbl3}\\
    \hline
    \rowcolor{gray!30} Наименование показателя & Значение \\ \hline
    \endfirsthead
    \caption*{\hspace{3pt}\emph{Продолжение таблицы \ref{sec1:tbl3}\hfill\vspace{-0.5\baselineskip}}}\\
    \hline
    \rowcolor{gray!30} Наименование показателя & Значение \\ \hline
    \endhead
    \hline
    \endfoot
    \endlastfoot
    \raggedright Электрическая прочность изоляции цепей с рабочим напряжением не более 60 В \\ по ГОСТ IEC 60255-5-2014, п. 6.1.4 & \centering\arraybackslash 500 В действующего значения, 50 Гц, 1 мин \\ \hline
    \raggedright Электрическая прочность изоляции цепей с рабочим напряжением более 60 В \\ по ГОСТ IEC 60255-5-2014, п. 6.1.4 & \centering\arraybackslash 2000 В действующего значения, 50 Гц, 1 мин \\ \hline
    \raggedright Импульсное напряжение 1,2/50 мкс для цепей с рабочим напряжением не более 60 В \\ по ГОСТ IEC 60255-5-2014, п. 6.1.3 & \centering\arraybackslash 1 кВ \\ \hline
    \raggedright Импульсное напряжение 1,2/50 мкс для цепей с рабочим напряжением более 60 В \\ по ГОСТ IEC 60255-5-2014, п. 6.1.3 & \centering\arraybackslash 5 кВ \\ \hline
  \end{longtable}
  \endgroup
\item При повторных испытаниях шкафа испытательное напряжение не должно превышать 85~\% от вышеуказанных значений.
\end{enumsubsub}

\subsubsection{Цепи оперативного питания}

\begin{enumsubsub}
\item \textls[-10]{Питание шкафа осуществляется от цепей постоянного оперативного тока номинальным напряжением 220~В. Параметры цепей питания ИЭУ в составе шкафа приведены в руководстве по эксплуатации \ManualUnitGeneral.}
\item Шкаф правильно функционирует при изменении напряжения оперативного тока в диапазоне от 0,8 до 1,1 номинального значения.
\item Контакты выходных реле шкафа не замыкаются ложно при подаче и снятии напряжения оперативного тока с перерывом любой длительности.
\end{enumsubsub}

\subsubsection{Характеристики аналоговых цепей}
\begin{enumsubsub}
\item \textls[-10]{Характеристики аналоговых цепей ИЭУ в составе шкафа приведены в руководстве по эксплуатации \ManualUnitGeneral.}
\end{enumsubsub}

\subsubsection{Характеристики дискретных входов и выходов}
\begin{enumsubsub}
\item Характеристики дискретных входов и выходов ИЭУ в составе шкафа приведены в руководстве по эксплуатации \ManualUnitGeneral.
\end{enumsubsub}

\subsubsection{Электромагнитная совместимость}
\begin{enumsubsub}
\item \textls[-3]{Шкаф соответствует требованиям по электромагнитной совместимости согласно ГОСТ IEC 61439-1. Сведения об устойчивости ИЭУ к воздействию помех приведены в руководстве по эксплуатации \ManualUnitGeneral.}
\end{enumsubsub}

\subsubsection{Надежность}

\begin{enumsubsub}
\item Средняя наработка на отказ шкафа составляет не менее 125 000~ч.
\item Среднее время восстановления работоспособного состояния ИЭУ в составе шкафа при наличии полного комплекта запасных блоков составляет не более 2~ч с учетом времени нахождения неисправности.
\item Средний срок службы шкафа -- не менее 25 лет при условии проведения требуемых технических мероприятий по обслуживанию с заменой, при необходимости, материалов и комплектующих, имеющих меньший срок службы.
\item Средний срок сохраняемости шкафа в упаковке поставщика составляет три года.
\item \textls[-5]{Периодичность планового технического обслуживания определяется характеристиками ИЭУ и комплектующих входящих в состав шкафа, а также моделью организации работ в эксплуатирующей организации:}
  \begin{itemize}
  \item по графику;
  \item по состоянию, с автоматизированным дистанционным мониторингом или без него.
  \end{itemize}
\item Минимальный срок периодического технического обслуживания шкафа должен составлять не менее 3 лет.
\end{enumsubsub}

\subsubsection{Защитное заземление}

\begin{enumsubsub}
\item Защитное заземление шкафа осуществляется через отдельный заземляющий проводник подключенный к специальному элементу заземления.
\item Шкаф имеет элементы для заземления (шина, болты, винты). Диаметр болтов (винтов, шпилек) для заземления, размеры контактной площадки, к которым присоединяются защитные провода заземления, должны соответствовать ГОСТ 12.2.007.0-75. Элементы заземления размещаются внутри шкафа.
\item Все металлические части устройств в составе шкафа и металлоконструкции, которые могут оказаться под напряжением вследствие нарушения изоляции, заземляются.
\item В соответствии с ГОСТ 12.2.007.0-75 обеспечивается непрерывность цепи защитного заземления. Значение сопротивления между заземляющим болтом для заземления шкафа и каждой доступной прикосновению металлической нетоковедущей частью оборудования не должно превышать 0,1~Ом (заземляющая цепь должна быть непрерывной).
\item Проводники защитного заземления отличаются от других цепей и имеют двухцветную изоляцию желто-зеленого цвета по всей своей длине. Главный элемент заземления шкафа имеет знак «заземления».
\end{enumsubsub}

% ============================================================
% 3. СОСТАВ ШКАФА
% ============================================================
{\color{unidarkgreen}\subsection{Состав шкафа}}

\subsubsection{Конструктивное исполнение}

\begin{enumsubsub}
\item Шкаф является защищённым низковольтным комплектным устройством (НКУ). Шкаф представляет собой металлоконструкцию из специализированного профиля с двухсторонним обслуживанием.
\IfSubStr{\twodoorcab}{\devicecode}{%
  \item Габариты шкафа 1600×600×2200~мм (Ш×Г×В, включая цоколь 200~мм).
}{\item Габариты шкафа 800×600×2200~мм (Ш×Г×В, включая цоколь 200~мм).}%
\item Габаритные и установочные размеры шкафа представлены на чертеже в приложении \ref{app:size}.
\item Несущей конструкцией шкафа является цельносварная рама. Для крепления изделия применяется монтажная конструкция в виде цоколя. Изделия крепятся к цоколю болтовыми соединениями.
\item Для крепления оборудования внутри шкафа предусмотрена монтажная панель.
\IfSubStr{\twodoorcab}{\devicecode}{%
  \item Оболочка шкафа предусматривает двери: передняя дверь (двухстворчатая -- сплошная обзорная) и задняя дверь (двухстворчатая -- глухая).
}{\item Оболочка шкафа предусматривает двери: передняя дверь (одностворчатая -- сплошная обзорная) и задняя дверь (двухстворчатая -- глухая).}%
\item Двери обеспечивают свободное открывание и закрывание на угол не менее 110°, задняя дверь оборудована упорами для фиксации двери в открытом положении. Двери шкафов снабжены замками, предотвращающими несанкционированный доступ.
\item На монтажной плите шкафа располагаются:
  \begin{itemize}
  \item табличка с диспетчерским/оперативным наименованием шкафа;
  \item шильдик шкафа;
  \item лампа сигнализации (HL);
  \item блок интерфейсный ИЧМ;
  \item оперативные переключатели (SA);
  \item прозрачный пластиковый карман;
  \item блоки испытательные (SG), через которые подключаются входные аналоговые цепи шкафа от трансформаторов тока и напряжения;
  \item технологическое отверстие.
  \end{itemize}
  Крепление интеллектуального электронного устройства предусмотрено на DIN-рейках внутри шкафа.
\item Подвод кабельных связей к шкафу осуществляется снизу через дно шкафа. Кабельные связи подключаются к клеммным рядам, расположенным на правой и левой боковинах. Доступ к клеммным рядам, устройствам и монтажу шкафа обеспечивается при открытии задней двери.
\item Монтаж аппаратов шкафа между собой выполнен медными проводами. Номинальное сечение проводов не менее 2,5~мм$^2$ для токовых цепей, не менее 0,75~мм$^2$ - для остальных цепей. Допускается отклонение от указанных требований при условии обеспечения выполнения требований к термической стойкости и механической прочности.
\item Допускается подключение к клеммным зажимам одного проводника сечением не более 10~мм$^2$ или двух проводников сечением не более 2,5~мм$^2$ для цепей тока. Для остальных цепей допускается подключение одного проводника сечением не более 6~мм$^2$ или двух проводников сечением не более 1,5~мм$^2$.
\end{enumsubsub}

\subsubsection{Схемы шкафа}
\begin{enumsubsub}
\item Чертеж общего вида представлен в приложении \ref{app:vo};
\item Схема электрическая принципиальная представлена в приложении \ref{app:a3};
\item Схема электрическая соединений рядов зажимов представлена в приложении \ref{app:a41};
\item Перечень элементов шкафов представлен в приложении \ref{app:spec}.
\end{enumsubsub}

% ============================================================
% 4. УСТРОЙСТВО И РАБОТА
% ============================================================
{\color{unidarkgreen}\subsection{Устройство и работа} \label{sec1:4}}

Принцип работы функций защит и автоматики ИЭУ, входящих в состав шкафа, описаны в \ManualUnitDefinite.

% ============================================================
% 5. СРЕДСТВА ИЗМЕРЕНИЙ, ИНСТРУМЕНТ И ПРИНАДЛЕЖНОСТИ
% ============================================================
{\color{unidarkgreen}\subsection{Средства измерений, инструмент и принадлежности}}

Перечень оборудования и средств измерений, необходимых для проведения эксплуатационных проверок шкафа, приведён в приложении \ref{app:unit}.

% ============================================================
% 6. МАРКИРОВКА И ПЛОМБИРОВАНИЕ
% ============================================================
{\color{unidarkgreen}\subsection{Маркировка и пломбирование}}

\begin{enumsub}
\item Шкаф имеет маркировку в соответствии с конструкторской документацией. Маркировка выполнена в соответствии с ГОСТ 18620-86 способом, обеспечивающим её чёткость и сохраняемость.
\item Информационная табличка (шильдик) размещается в правом верхнем углу на передней двери шкафа и дублируется на монтажной панели шкафа с лицевой стороны (см. рисунок \ref{section1:fig_1}).
  \medskip
  \begin{figure}[H]
    \centering
    \includegraphics[width=0.6\textwidth,height=0.6\textheight,keepaspectratio]{\genpath/img/Nameplate.png}
    \captionof{figure}{Эскизный чертеж информационной таблички типового шкафа \\
      1 – код шкафа; 2 – дата выпуска шкафа в формате ММ.ГГГГ; 3 – поле для QR-кода}
    \label{section1:fig_1}
  \end{figure}
  \medskip
\item На табличке (шильдике) в дополнение к текстовой информации наносится QR-код, содержащий:
  \begin{itemize}
  \item наименование шкафа;
  \item шифр шкафа;
  \item основные функции ИЭУ шкафа;
  \item номинальный вторичный переменный ток;
  \item номинальная частота;
  \item номинальное переменное напряжение;
  \item напряжение оперативного постоянного тока;
  \item дата (месяц, год) выпуска шкафа в формате ММ.ГГГГ.
  \end{itemize}
\item Все элементы шкафа имеют обозначение, соответствующее схеме шкафа и состоящее из буквенного обозначения и порядкового номера, проставленного после буквенного обозначения (например, SG1).
\item Транспортная маркировка тары выполнена в соответствии с ГОСТ 14192-96. Манипуляционные знаки наносят в верхней части двух вертикальных смежных сторон.
\end{enumsub}

% ============================================================
% 7. УПАКОВКА
% ============================================================
{\color{unidarkgreen}\subsection{Упаковка}}

\begin{enumsub}
\item Упаковка шкафа проводится по конструкторской документации предприятия-изготовителя для условий хранения, транспортирования и допустимых сроков сохраняемости.
\item В соответствии с ГОСТ 23216-78 для изделий применяется транспортная тара двух видов:
  \begin{itemize}
  \item ТЭ: ящики дощатые;
  \item ТФ: ящики фанерные и из древесно-волокнистых плит.
  \end{itemize}
\end{enumsub}